\documentclass[a4paper,12pt]{article}
\usepackage{cmap}
% --- Кодировки и шрифты ---
\usepackage[utf8]{inputenc}
\usepackage[T1, T2A]{fontenc} 
\usepackage[russian]{babel}

% --- Математические пакеты ---
\usepackage{amsmath}
\usepackage{amssymb}
\usepackage{amsthm}
\usepackage{mathtools}
\usepackage{tikz}
\usepackage{mathrsfs}

% --- Настройка полей ---
\usepackage[left=2cm, right=2cm, top=2cm, bottom=2cm]{geometry}

% --- Гиперссылки ---
\usepackage[colorlinks=true, linkcolor=blue, urlcolor=blue]{hyperref}

% --- ДОБАВЛЕННЫЕ ПАКЕТЫ ДЛЯ ОФОРМЛЕНИЯ ---
\usepackage[dvipsnames]{xcolor}
\usepackage[most]{tcolorbox}

% --- КОЛОНТИТУЛЫ (НАВИГАЦИЯ) ---
\usepackage{fancyhdr}
\pagestyle{fancy}
\fancyhf{}
\setlength{\headheight}{15pt}

% Команда для обновления кнопок
\newcommand{\setupnav}[2]{
	\fancyhead[L]{\hyperlink{toc}{\textbf{[Меню]}}}
	\fancyhead[R]{
		\ifx&#1& \textcolor{gray}{[$\leftarrow$ Назад]} \else \hyperlink{bilet#1}{[$\leftarrow$ Назад]} \fi
		\quad | \quad
		\ifx&#2& \textcolor{gray}{[Вперед $\rightarrow$]} \else \hyperlink{bilet#2}{[Вперед $\rightarrow$]} \fi
	}
}
\fancyfoot[C]{\thepage}

% --- ЦВЕТА ---
\definecolor{defmyorange}{RGB}{255, 140, 0}
\definecolor{thmblue}{RGB}{0, 102, 204}
\definecolor{proofpurp}{RGB}{128, 0, 128}
\definecolor{highlighter}{RGB}{255, 215, 0}
\definecolor{exgray}{RGB}{100, 100, 100}

% --- Определение окружений ---
\theoremstyle{plain}
\newtheorem{theorem}{Теорема}[section]
\newtheorem{lemma}[theorem]{Лемма}
\newtheorem{proposition}[theorem]{Утверждение}
\newtheorem{corollary}[theorem]{Следствие}

\theoremstyle{definition}
\newtheorem{definition}{Определение}[section]
\newtheorem{example}{Пример}[section]

\theoremstyle{remark}
\newtheorem{remark}{Замечание}[section]
\newtheorem*{idea}{Идея доказательства}

% --- МАГИЯ ПЕРЕОПРЕДЕЛЕНИЯ ---
\renewenvironment{definition}[1][]{
	\begin{tcolorbox}[enhanced, colback=defmyorange!5, colframe=defmyorange, coltitle=white, fonttitle=\bfseries, title={Определение: #1}, attach boxed title to top left={yshift=-2mm, xshift=2mm}, boxed title style={colback=defmyorange, sharp corners}, sharp corners, top=12pt, breakable]
	}{\end{tcolorbox}}

\renewenvironment{theorem}[1][]{
	\begin{tcolorbox}[enhanced, colback=thmblue!5, colframe=thmblue, fonttitle=\bfseries, title={Теорема #1}, sharp corners=south, drop shadow, breakable]
	}{\end{tcolorbox}}

\renewenvironment{lemma}[1][]{
	\begin{tcolorbox}[enhanced, colback=thmblue!5, colframe=thmblue, fonttitle=\bfseries, title={Лемма #1}, sharp corners=south, breakable]
	}{\end{tcolorbox}}

\renewenvironment{proposition}[1][]{
	\begin{tcolorbox}[enhanced, colback=thmblue!5, colframe=thmblue, fonttitle=\bfseries, title={Утверждение #1}, sharp corners=south, breakable]
	}{\end{tcolorbox}}

\renewenvironment{corollary}[1][]{
	\begin{tcolorbox}[enhanced, colback=thmblue!2, colframe=thmblue!70, fonttitle=\bfseries, title={Следствие #1}, sharp corners, breakable]
	}{\end{tcolorbox}}

\renewenvironment{proof}{
	\begin{tcolorbox}[
		enhanced, 
		breakable, 
		colback=proofpurp!2,    % Светлый фон
		colframe=proofpurp,     % Цвет рамки
		coltitle=white,         % Цвет текста в заголовке
		fonttitle=\bfseries, 
		title={Доказательство}, % Заголовок в рамке
		sharp corners,          % Острые углы как у теорем
		boxrule=1pt,            % Толщина рамки
		toptitle=2mm,           % Отступы в заголовке
		bottomtitle=2mm
		]
	}{\qed \end{tcolorbox}}

\renewenvironment{example}[1][]{
	\begin{tcolorbox}[enhanced, colback=exgray!5, colframe=exgray, fonttitle=\bfseries, title={Пример #1}, sharp corners, boxrule=0.5pt, breakable]
	}{\end{tcolorbox}}

\renewenvironment{remark}[1][]{
	\vspace{2mm}\noindent\textit{Замечание #1.}~
}{ \vspace{2mm} }

\setcounter{tocdepth}{2}


% --- НОВЫЙ ЦВЕТ ДЛЯ ИДЕИ ---
\definecolor{ideared}{RGB}{196, 30, 58} % Красивый красный (кардинал)

% --- МАГИЯ ПЕРЕОПРЕДЕЛЕНИЯ IDEA ---
\renewenvironment{idea}[1][]{
	\begin{tcolorbox}[
		enhanced,
		breakable,
		colback=ideared!5, 
		colframe=ideared, 
		borderline west={2pt}{0pt}{ideared},
		coltitle=ideared,
		fonttitle=\bfseries,
		title={Идея доказательства: #1},
		attach title to upper,
		after title={\par\nobreak\vspace{1mm}},
		sharp corners
		]
	}{\end{tcolorbox}}


% --- Заголовок ---
\title{Билеты по функциональному анализу}
\author{}
\date{\today}

\begin{document}
	
	\maketitle
	\hypertarget{toc}{} 
	
	\begin{abstract}
		Билеты по функциональному анализу для подготовки к экзамену. Оформлено по разделам программы (2026).
	\end{abstract}
	
	\tableofcontents
	\newpage
	
	% I. МЕТРИЧЕСКИЕ ПРОСТРАНСТВА
	\section{Метрические пространства}
	\hypertarget{bilet1}{}\setupnav{}{2} % --- Содержимое файла materials/1_1.tex ---

\subsection{Метрические и топологические пространства. Примеры: $l_p, C_p[a, b] (p \geqslant 1), C^n[a, b]$. Неравенства Гёльдера и Минковского}

\begin{definition}[Метрическое пространство]
	Метрическим пространством называется множество $X$ с метрикой $\rho: X^2 \to \mathbb{R}$, обладающей следующими свойствами:
	\begin{enumerate}
		\item $\forall x, y \in X : \rho(x, y) \geqslant 0$, причем $\rho(x, y) = 0 \Leftrightarrow x = y$
		\item $\forall x, y \in X : \rho(x, y) = \rho(y, x)$
		\item $\forall x, y, z \in X : \rho(x, z) \leqslant \rho(x, y) + \rho(y, z)$ (неравенство треугольника)
	\end{enumerate}
\end{definition}

\begin{definition}[Топологическое пространство]
	Топологическим пространством называется множество $X$ с системой $\tau \subset 2^X$, обладающей следующими свойствами:
	\begin{enumerate}
		\item $\varnothing, X \in \tau$
		\item $\forall i \in \{1, 2, \dots, n\}, G_i \in \tau$ верно $\bigcap_{i=1}^n G_i \in \tau$
		\item $\forall A$ — индексирующего множества, $G_{\alpha} \in \tau, \forall \alpha \in A$ верно $\bigcup_{\alpha \in A} G_{\alpha} \in \tau$
	\end{enumerate}
	Система $\tau$ называется \textbf{топологией} на множестве $X$, а элементы системы $\tau$ — \textbf{открытыми множествами}.
\end{definition}

\begin{theorem}[Неравенства Гёльдера и Минковского]
	Пусть $E$ — измеримое. Пусть $f, g : E \to [0, +\infty]$ измеримы, и $1 < p < \infty, \frac{1}{p} + \frac{1}{q} = 1$. Тогда справедливо \textbf{неравенство Гёльдера}
	\[ \int_E fg d\mu \leqslant \left( \int_E f^p d\mu \right)^{\frac{1}{p}} \cdot \left( \int_E g^q d\mu \right)^{\frac{1}{q}} \]
	и \textbf{неравенство Минковского}
	\[ \left( \int_E (f + g)^p d\mu \right)^{\frac{1}{p}} \leqslant \left( \int_E f^p d\mu \right)^{\frac{1}{p}} + \left( \int_E g^q d\mu \right)^{\frac{1}{p}} \]
\end{theorem}

\begin{remark}
	Таблица метрических пространств и их свойства.
\end{remark}

\begin{definition}[Подпространство]
	Пусть $X$ — МП, $Y \subset X$. \textbf{Подпространством} пространства $X$ называется метрическое пространство $Y$ с метрикой, являющейся сужением метрики на $X$.
\end{definition}

\begin{definition}[Ограниченность]
	Пусть $X$ — МП. Множество $Y \subset X$ называется \textbf{ограниченным}, если выполнено условие $\sup_{x,y \in Y} \rho(x, y) < +\infty$.
\end{definition}

\begin{remark}
	Далее в этом разделе мы стремимся показать, что каждое МП является ТП с топологией определенного вида, заданной метрикой пространства.
\end{remark}

\begin{definition}[Шары]
	Пусть $X$ — МП, $x \in X, r > 0$.
	\begin{itemize}
		\item \textbf{Открытым шаром} называется множество $B(x, r) := \{y \in X : \rho(y, x) < r\}$
		\item \textbf{Замкнутым шаром} называется множество $\overline{B}(x, r) := \{y \in X : \rho(y, x) \leqslant r\}$
	\end{itemize}
\end{definition}

\begin{remark}
	Также достаточно часто используются альтернативные обозначения шаров: $B_r(x) := B(x, r)$ и $\overline{B}_r(x) := \overline{B}(x, r)$.
\end{remark}

\begin{definition}[Внутренность]
	Пусть $X$ — МП, $M \subset X$. Точка $x \in X$ называется \textbf{внутренней точкой} множества $M$, если существует $r > 0$ такое, что $B(x, r) \subset M$. \textbf{Внутренностью} множества $M$ называется множество $\text{int } M$ всех его внутренних точек. Множество $M$ называется \textbf{открытым}, если $\text{int } M = M$.
\end{definition}

\begin{definition}[Замыкание]
	Пусть $(X, \rho)$ — МП, $M \subset X$. Точка $x \in X$ называется \textbf{точкой прикосновения} множества $M$, если для любого $r > 0$ выполнено условие $B(x, r) \cap M \neq \varnothing$. \textbf{Замыканием} множества $M$ называется множество $\overline{M}$ всех его точек прикосновения. Множество $M$ называется \textbf{замкнутым}, если $\overline{M} = M$.
\end{definition}

\begin{theorem}
	Пусть $X$ — МП, $M \subset X$. Тогда множество $M$ открыто $\Leftrightarrow$ множество $X \setminus M$ замкнуто.
\end{theorem}

\begin{theorem}[Двойственность замкнутых и открыхы]
	Пусть $X$ — МП. Тогда:
	\begin{enumerate}
		\item Для любой системы открытых множеств $\{G_\alpha\}_{\alpha \in A} \subset 2^X$ множество $\bigcup_{\alpha \in A} G_\alpha$ тоже является открытым
		\item Для любых открытых множеств $G_1, G_2 \subset X$ множество $G_1 \cap G_2$ тоже является открытым
		\item Для любой системы замкнутых множеств $\{F_\alpha\}_{\alpha \in A} \subset 2^X$ множество $\bigcap_{\alpha \in A} F_\alpha$ тоже является замкнутым
		\item Для любых замкнутых множеств $F_1, F_2 \subset X$ множество $F_1 \cup F_2$ тоже является замкнутым
	\end{enumerate}

\end{theorem}

\begin{theorem}
	Пусть $X$ — МП. Тогда:
	\begin{enumerate}
		\item Для любого $x \in X$ и $r > 0$ множество $B(x, r)$ — открытое
		\item Для любого $x \in X$ и $r > 0$ множество $\overline{B}(x, r)$ — замкнутое
		\item Для любого множества $M \subset X$ множество $\text{int } M$ — открытое, причем наибольшее по включению открытое множество, содержащееся в $M$
		\item Для любого множества $M \subset X$ множество $\overline{M}$ — замкнутое, причем наименьшее по включению замкнутое множество, содержащее $M$
	\end{enumerate}
\end{theorem}

\begin{definition}[Плотность]
	Пусть $X$ — МП. Множество $A \subset X$ называется:
	\begin{itemize}
		\item \textbf{Плотным} в множестве $B \subset X$, если $B \subset \overline{A}$
		\item \textbf{Всюду плотным}, если $X = \overline{A}$
	\end{itemize}
\end{definition}

\begin{definition}[Всюду плотное множество]
	Пусть $(X, \rho)$ "--- метрическое пространство. Подмножество $S \subset X$ называется \textbf{всюду плотным} в $X$, если любой элемент пространства $X$ можно сколь угодно точно приблизить элементами из $S$.
	
	Формально: для любой точки $x \in X$ и любой погрешности $\varepsilon > 0$ найдется элемент $s \in S$, лежащий на расстоянии меньше $\varepsilon$ от $x$:
	\[ \forall x \in X, \; \forall \varepsilon > 0 \quad \exists s \in S : \rho(x, s) < \varepsilon \]
\end{definition}

\begin{definition}[Сепарабельность]
	МП $X$ называется \textbf{сепарабельным}, если в $X$ существует не более чем счетное всюду плотное множество.
\end{definition}

\begin{definition}[Сходимость]
	Пусть $X$ — МП. Последовательность $\{x_n\} \subset X$ \textbf{сходится} к точке $x \in X$, если $\rho(x_n, x) \to 0$ при $n \to \infty$. Обозначение — $x_n \xrightarrow[X]{} x$.
\end{definition}

\begin{definition}[Непрерывность в МП]
	Пусть $X, Y$ — МП, $f: X \to Y$. Отображение $f$ называется \textbf{непрерывным в точке} $x \in X$, если выполнено одно из следующих условий:
	\begin{enumerate}
		\item Для любого $\varepsilon > 0$ существует $\delta > 0$ такое, что $f(B(x, \delta)) \subset B(f(x), \varepsilon)$
		\item Для любой $\{x_n\} \subset X$ такой, что $x_n \to x$, выполнено $f(x_n) \to f(x)$
	\end{enumerate}
	Отображение $f$ называется \textbf{непрерывным}, если оно непрерывно во всех точках из $X$.
\end{definition}

\begin{remark}
	Для ТП непрерывность определяется похожим образом и согласуется с непрерывностью в МП. Пусть $X, Y$ — ТП, $f: X \to Y$. Отображение $f$ называется \textbf{непрерывным в точке} $x \in X$, если выполнено:
	\[ \forall V(f(x)) \in \tau_Y \exists U(x) \in \tau_X : f(U(x)) \subset V(f(x)). \]
\end{remark}

\begin{remark}
	Можно заметить, что если прообраз любого открытого множества будет открыт, то и прообраз любого замкнутого будет замкнут.
\end{remark}

\begin{definition}[Гомеоморфизм]
	Если отображение между МП $f: X \to Y$ взаимнооднозначно и взаимнонепрерывно, то есть $f^{-1}$ и $f$ непрерывны, то $f$ — \textbf{гомеоморфизм}. Пространства $X$ и $Y$ называются гомеоморфными.
\end{definition}

\begin{remark}
	В случае топологических пространств, гомеоморфизм это биекция между $X$ и $Y$, а также между топологиями на них.
\end{remark}

\begin{definition}[Изометрия]
	Если гомеоморфизм $f$ сохраняет расстояния:
	\[ \rho(x_1, x_2) = \rho'(f(x_1), f(x_2)), \]
	то он называется \textbf{изометрией}. Пространства $X$ и $Y$ называются изометричными.
\end{definition}

\begin{example}[Пространство $l_p$, $p \geqslant 1$]
	Это пространство числовых последовательностей $x = (x_1, x_2, \dots)$, для которых сходится ряд из модулей в степени $p$:
	\[ l_p = \left\{ x = (x_k)_{k=1}^\infty \mathrel{\bigg|} \sum_{k=1}^\infty |x_k|^p < \infty \right\} \]
	Метрика задается формулой:
	\[ \rho(x, y) = \left( \sum_{k=1}^\infty |x_k - y_k|^p \right)^{1/p} \]
	\textbf{Почему это метрика?} Аксиома треугольника для этой метрики называется \textit{неравенством Минковского для сумм}:
	\[ \left( \sum |x_k + y_k|^p \right)^{1/p} \leqslant \left( \sum |x_k|^p \right)^{1/p} + \left( \sum |y_k|^p \right)^{1/p} \]
\end{example}

\begin{example}[Пространство $C_p{[a, b]}$, $p \geqslant 1$]
	Множество элементов — это \textbf{непрерывные} функции на отрезке $[a, b]$ (то же множество, что и в $C[a, b]$).
	Однако метрика задается не через максимум, а через интеграл:
	\[ \rho(f, g) = \left( \int_a^b |f(t) - g(t)|^p \, dt \right)^{1/p} \]
	\textbf{Нюанс:} В этой метрике последовательность функций может сходиться, даже если она не сходится равномерно (пики могут становиться очень узкими, но высокими, главное чтобы площадь стремилась к нулю).
	\\
	Аксиома треугольника здесь обеспечивается \textit{неравенством Минковского для интегралов}.
\end{example}



\begin{example}[Пространство $C^n{[a, b]}$]
	Это пространство функций, имеющих на отрезке $[a, b]$ непрерывные производные до порядка $n$ включительно.
	Метрика учитывает близость не только самих функций, но и всех их производных:
	\[ \rho(f, g) = \sum_{k=0}^n \max_{t \in [a, b]} |f^{(k)}(t) - g^{(k)}(t)| \]
	Или эквивалентная ей:
	\[ \rho(f, g) = \max_{k \in \{0, \dots, n\}} \left( \max_{t \in [a, b]} |f^{(k)}(t) - g^{(k)}(t)| \right) \]
	Сходимость в этом пространстве означает равномерную сходимость самих функций и всех их производных до $n$-го порядка.
\end{example}\newpage
	
	% II. ПОЛНЫЕ МЕТРИЧЕСКИЕ ПРОСТРАНСТВА
	\section{Полные метрические пространства}
	\hypertarget{bilet2}{}\setupnav{1}{3} % --- Содержимое файла materials/2_1.tex ---
\subsection{Теорема о вложенных шарах и теорема Бэра}
\begin{definition}[Полное метрическое пространство]
	МП \textit{полное}, если любая фундаментальная последовательность сходится.
\end{definition}

\begin{theorem}[О вложенных шарах]\label{kantor-compact}
	\textbf{Условия:}
	\begin{itemize}
		\item $X$ --- полное метрическое пространство;
		\item $\{\overline{B}_{r_n}(x_n)\}$ --- последовательность вложенных замкнутых шаров ($\overline{B}_{r_{n+1}} \subset \overline{B}_{r_n}$);
		\item $r_n \to 0$ при $n \to \infty$.
	\end{itemize}
	\textbf{Следовательно:}
	\begin{itemize}
		\item $\bigcap\limits_{n \in \mathbb{N}} \overline{B}_{r_n}(x_n) \neq \varnothing$;
		\item Пересечение состоит ровно из одной точки: $\bigcap\limits_{n \in \mathbb{N}} \overline B_{r_n}(x_n) = \{x^*\}$.
	\end{itemize}
\end{theorem}

\begin{proof}
	\begin{enumerate}
		\item \textbf{Фундаментальность:} В силу вложенности шаров $\overline{B}_{r_{n+p}} \subset \overline{B}_{r_n}$ расстояние между центрами $\rho(x_n, x_{n+p}) \leqslant r_n$. Так как по условию $r_n \to 0$, последовательность $\{x_n\}$ является фундаментальной.
		\item \textbf{Сходимость:} Поскольку пространство $X$ полно по определению, существует точка $x^* \in X$, такая что $x_n \to x^*$.
		\item \textbf{Принадлежность шарам:} Для любого фиксированного $N$ все точки $\{x_n\}$ при $n \geqslant N$ лежат в шаре $\overline B_{r_N}(x_N)$. Так как шар является замкнутым множеством, он содержит и свой предел: $x^* \in \overline B_{r_N}(x_N)$. Это верно для любого $N$, следовательно, $x^* \in \bigcap_{n \in \mathbb{N}} \overline B_{r_n}(x_n)$.
		\item \textbf{Единственность:} Если бы существовала вторая точка $y^* \in \bigcap \overline B_{r_n}(x_n)$, то расстояние $\rho(x^*, y^*) \leqslant 2r_n$ для любого $n$. Устремляя $n \to \infty$, получаем $\rho(x^*, y^*) = 0 \Rightarrow x^* = y^*$.
	\end{enumerate}
\end{proof}

\begin{definition}[Нигде не плотное множество]
	Множество $M \subset X$ называется \textit{нигде не плотным} в $X$, если внутренность его замыкания пуста:
	\[ \text{int}(\overline{M}) = \varnothing \]
\end{definition}

\begin{definition}[Критерий «дырявости»]
	Эквивалентно, $M$ \textit{нигде не плотно}, если в любом открытом шаре $B \subset X$ содержится подшар $B_1 \subset B$, не пересекающийся с $M$:
	\[ B_1 \cap M = \varnothing \]
\end{definition}

\begin{definition}[Всюду плотное множество]
	Множество $M \subset X$ называется \textit{всюду плотным} в $X$, если его замыкание совпадает со всем пространством:
	\[ \overline{M} = X \]
	(То есть в любой окрестности любой точки из $X$ есть хотя бы одна точка из $M$).
\end{definition}

\begin{theorem}[Бэра]\label{Ber}
	\textbf{Условия:}
	\begin{itemize}
		\item $X$ --- полное метрическое пространство;
		\item $M_n \subset X$ --- последовательность множеств, нигде не плотных в $X$.
	\end{itemize}
	\textbf{Следовательно:}
	\begin{itemize}
		\item Пространство $X$ нельзя представить в виде счетного объединения таких множеств: $X \neq \bigcup\limits_{n \in \mathbb{N}} M_n$.
	\end{itemize}
\end{theorem}

\begin{proof}
	Предположим от противного: $X = \bigcup_{n=1}^{\infty} M_n$.
	\begin{enumerate}
		\item \textbf{Построение:} 
		Возьмем произвольный шар $\overline{B}_0$. 
		Так как $M_1$ нигде не плотно, внутри открытого шара $B_0$ существует замкнутый подшар $\overline{B}_1$, не пересекающийся с $M_1$. Радиус выберем $r_1 \leqslant 1$.
		
		\item \textbf{Индукция:} 
		Для каждого $n$ внутри открытого шара $B_{n-1}$ найдем замкнутый подшар $\overline{B}_n$ такой, что $\overline{B}_n \cap M_n = \varnothing$ и $r_n \leqslant \frac{1}{2^n}$. (сначала определение, а радиус можем любой взять) Возможность такого выбора гарантирована определением нигде не плотного множества.
		
		\item \textbf{Предел:} 
		Мы получили систему вложенных замкнутых шаров $\overline{B}_n$, радиусы которых стремятся к нулю. В силу \textit{полноты} пространства $X$, существует точка $x^* \in \bigcap_{n=1}^{\infty} \overline{B}_n$.
		
		\item \textbf{Противоречие:} 
		Поскольку $x^* \in \overline{B}_n$ для любого $n$, то $x^* \notin M_n$ ни для одного $n$ (так как $\overline{B}_n \cap M_n = \varnothing$). Следовательно, $x^* \notin \bigcup M_n$. 
		Но по предположению $X = \bigcup M_n$, а значит $x^*$ обязана была попасть в объединение. Противоречие.
	\end{enumerate}
\end{proof}

\begin{idea}{Ловушка для понимания (Почему $\overline{M}_n$?)}
	\textbf{Важный нюанс:} В доказательстве мы строим шары так, чтобы они не пересекались именно с \textit{замыканием} $\overline{M}_n$. 
	
	\textbf{Почему это критично?} Если «уворачиваться» только от самого $M_n$, то предел последовательности шаров $x^*$ может оказаться предельной точкой для $M_n$ (лежать на его границе). Если граница принадлежит $M_n$, то $x^* \in M_n$, и противоречие исчезнет. 
	
	Выбирая шар $\overline{B}_n \cap \overline{M}_n = \varnothing$, мы создаем «зону безопасности», гарантирующую, что $x^*$ не только не совпадет с точками «пыли», но и не окажется на её «краю».
\end{idea}
\newpage
	\hypertarget{bilet3}{}\setupnav{2}{4} % --- Содержимое файла materials/2_2.tex ---
\subsection{Теорема Банаха о сжимающих отображениях}
\begin{theorem}[О сжимающем отображении (Принцип Банаха)]\label{banach-compact}
	\textbf{Условия:}
	\begin{itemize}
		\item $X$ --- полное метрическое пространство;
		\item $F: X \to X$ --- сжимающее отображение, т.е. существует такое $\lambda \in (0, 1)$, что для всех $x, y \in X$ выполнено:
		\[ \rho(F(x), F(y)) \leqslant \lambda \rho(x, y) \]
	\end{itemize}
	\textbf{Следовательно:}
	\begin{itemize}
		\item У отображения $F$ существует единственная неподвижная точка $x^* \in X$, такая что $F(x^*) = x^*$.
	\end{itemize}
\end{theorem}
 
\begin{proof}
	\begin{enumerate}
		\item \textbf{Построение последовательности (метод итераций):} Зафиксируем произвольное $x_0 \in X$ и определим последовательность итераций: $x_1 = F(x_0)$, $x_2 = F(x_1)$, ..., $x_{n+1} = F(x_n)$.
		
		\item \textbf{Оценка расстояния между соседними шагами:} Заметим, что по определению сжатия:
		\[ \rho(x_{k+1}, x_k) = \rho(F(x_k), F(x_{k-1})) \leqslant \lambda \rho(x_k, x_{k-1}) \leqslant \dots \leqslant \lambda^k \rho(x_1, x_0) \]
		
		\item \textbf{Фундаментальность:} Оценим расстояние между $x_n$ и $x_{n+p}$, используя неравенство треугольника и сумму геометрической прогрессии:
		\[ \rho(x_{n+p}, x_n) \leqslant \sum_{i=n}^{n+p-1} \rho(x_{i+1}, x_i) \leqslant \sum_{i=n}^{n+p-1} \lambda^i \rho(x_1, x_0) \leqslant \frac{\lambda^n}{1 - \lambda} \rho(x_1, x_0) \]
		Так как $\lambda < 1$, при $n \to \infty$ величина $\lambda^n \to 0$, следовательно, последовательность $\{x_n\}$ фундаментальна.
		
		\item \textbf{Существование предела:} Поскольку $X$ полно, существует предел $x^* = \lim_{n \to \infty} x_n$. Переходя к пределу в равенстве $x_{n+1} = F(x_n)$ (учитывая, что любое сжатие автоматически непрерывно), получаем:
		\[ \lim_{n \to \infty} x_{n+1} = F(\lim_{n \to \infty} x_n) \Rightarrow x^* = F(x^*) \]
		
		\item \textbf{Единственность:} Допустим, есть вторая точка $y^* = F(y^*)$. Тогда:
		\[ \rho(x^*, y^*) = \rho(F(x^*), F(y^*)) \leqslant \lambda \rho(x^*, y^*) \]
		Так как $\lambda < 1$, это неравенство возможно только при $\rho(x^*, y^*) = 0$, то есть $x^* = y^*$.
	\end{enumerate}
\end{proof}\newpage
	
	% III. КОМПАКТНЫЕ МЕТРИЧЕСКИЕ ПРОСТРАНСТВА
	\section{Компактные метрические пространства}
	\hypertarget{bilet4}{}\setupnav{3}{5} \subsection{компактность и центрированная систем замкнутных множетсв}
\begin{definition}[Компактное пространство]
	Метрическое пространство $X$ называется \textit{компактным}, если для любой системы открытых множеств $\{G_\alpha\}_{\alpha \in A} \subset 2^X$ такой, что $\bigcup\limits_{\alpha \in A} G_\alpha = X$, существует конечный набор $\alpha_1, \dotsc, \alpha_n \in A$ такой, что $\bigcup\limits_{k = 1}^n G_{\alpha_k} = X$.
\end{definition}

\begin{remark}
	Определение выше можно сформулировать так: из любого открытого покрытия пространства $X$ можно выделить конечное подпокрытие. Или, эквивалентно, если система открытых множеств не содержит конечного покрытия пространства $X$, то она не является покрытием пространства $X$.
\end{remark}

\begin{example}
	Любое замкнутое и ограниченное множество в $\mathbb{R}^n$ будет компактным (в индуцированной метрике).
\end{example}

\begin{definition}[Центрированная система]\label{center-system}
	Пусть $X$ --- метрическое пространство. Система множеств $\{B_\alpha\}_{\alpha \in A} \subset 2^X$ называется \textit{центрированной}, если для любого конечного набора индексов $\alpha_1, \dotsc, \alpha_n \in A$ выполнено:
	\[ \bigcap\limits_{k = 1}^n B_{\alpha_k} \neq \varnothing \]
\end{definition}



\begin{theorem}[Критерий компактности в терминах замкнутых множеств]\label{thm3.1}
	Следующие условия эквивалентны:
	\begin{itemize}
		\item Метрическое пространство $X$ компактно;
		\item Любая центрированная система замкнутых множеств в $X$ имеет непустое пересечение ($\bigcap_{\alpha \in A} B_\alpha \neq \varnothing$).
	\end{itemize}
\end{theorem}\newpage % Был 5-й, стал 4-й
	\hypertarget{bilet5}{}\setupnav{4}{6} \subsection{Критерий компактности}
\begin{definition}[$\varepsilon$-сеть]\label{varepsilon-net}
	Пусть $M$~--- некоторое множество в метрическом пространстве $R$. Тогда множество $A$ из $R$ называется \textbf{$\varepsilon$-сетью} для $M$, если для любой точки $x \in M$ найдется хотя бы одна точка $a \in A$, что
	\[\rho(x, a) \leqslant \varepsilon.\]
\end{definition}

\begin{definition}[Ограниченность]
	Множество $M$ в метрическом пространстве $R$ называется \textbf{ограниченным}, если существует $B_\varepsilon(x_0)$ содержащий $M$.
\end{definition}

\begin{definition}[Вполне ограниченность, полная ограниченность]
	Множество $M$ в метрическом пространстве $R$ называется \textbf{вполне ограниченным}, если для него при любом $\varepsilon > 0$ существует \textbf{конечная $\varepsilon$-сеть}.
\end{definition}

\begin{remark}
	Если множество вполне ограничено, то его замыкание также вполне ограничено.
\end{remark}

\begin{remark}
	Из определения следует, что если само пространство вполне ограничено, то оно сепарабельно.
\end{remark}

\begin{remark}
	Из вполне ограниченности следует ограниченность.
\end{remark}

\begin{proof}
	Из вполне ограниченности ограниченность получается как объединение конечного числа ограниченных множеств.
\end{proof}

\begin{remark}
	В $n$-мерном евклидовом пространстве вполне ограниченность совпадает с обычной ограниченностью.
\end{remark}

\begin{theorem}[Критерий компактности]\label{thm3.2}
	Пусть $X$ --- метрическое пространство. Тогда следующие условия эквивалентны:
	\begin{enumerate}
		\item $X$ компактно;
		\item $X$ полно и вполне ограниченно;
		\item Из любой последовательности $\{x_n\} \subset X$ можно выделить сходящуюся подпоследовательность $\{x_{n_k}\}$ (т.е. $X$ — \textit{секвенциально компактно});
		\item Любое бесконечное множество $M \subset X$ имеет предельную точку.
	\end{enumerate}
\end{theorem}

\begin{proof}
	\begin{itemize}
		\item[\textbf{(1) $\Rightarrow$ (2)}]
		\begin{enumerate}
			\item \textbf{Вполне ограниченность.} Для любого $\varepsilon > 0$ рассмотрим покрытие $\{B(x, \varepsilon)\}_{x \in X}$. В силу компактности $X$, выделим конечное подпокрытие. Центры этих шаров образуют конечную $\varepsilon$-сеть.
			
			
			
			\item \textbf{Полнота.} Пусть $\{x_n\}$ — фундаментальная последовательность. Рассмотрим множества $F_n = \overline{\{x_k\}_{k \geqslant n}}$.
			\begin{itemize}
				\item Система $\{F_n\}$ — центрированная система замкнутых множеств (пересечение любого набора содержит хвост последовательности).
				\item По Теореме~\ref{thm3.1}, $\bigcap_{n} F_n \neq \varnothing$. Пусть $x_0 \in \bigcap F_n$.
				\item Так как диам($F_n$) $\to 0$ (в силу фундаментальности), то $x_0$ единственна и $x_n \to x_0$.
			\end{itemize}
		\end{enumerate}
		
		\item[\textbf{(2) $\Rightarrow$ (3)}] Пусть $\{x_n\} \subset X$.
		\begin{enumerate}
			\item \textbf{Построение.} Так как $X$ вполне ограничено, его можно покрыть конечным числом шаров радиуса $1$. В одном из них ($B_1$) содержится бесконечно много членов $\{x_n\}$.
			\item Далее покрываем $B_1$ шарами радиуса $1/2$. В одном из них ($B_2 \subset B_1$) снова бесконечно много членов.
			\item Продолжая процесс, получим вложенные шары $B_1 \supset B_2 \supset \dots$ с радиусами $r_k \to 0$, содержащие подпоследовательность $\{x_{n_k}\}$ (выбираем $x_{n_k} \in B_k$).
			\item \textbf{Сходимость.} $\{x_{n_k}\}$ фундаментальна (так как лежит в шарах стягивающегося радиуса). В силу полноты $X$, она сходится.
		\end{enumerate}
		
		\item[\textbf{(3) $\Rightarrow$ (1)}]
		\begin{enumerate}
			\item \textbf{Вполне ограниченность.} От противного: если конечной сети нет, то можно построить последовательность «растолканных» точек: $x_1$, затем $x_2$ на расстоянии $\rho > \varepsilon_0$ от первой, $x_3$ на расстоянии $\rho > \varepsilon_0$ от первых двух и т.д. У такой «редкой» последовательности нет сходящейся подпоследовательности. Противоречие.
			
			\item \textbf{Компактность.} От противного: пусть покрытие $\{G_\alpha\}$ не имеет конечного подпокрытия. 
			Будем уменьшать масштаб:
			\begin{itemize}
				\item Разделим всё пространство на конечное число мелких шаров. Хотя бы один из них ($B_1$) — «плохой» (не покрывается конечно).
				\item Разделим этот «плохой» шар на еще более мелкие. Снова найдем «плохой» подшар $B_2 \subset B_1$.
				\item Повторяя процесс, получим вложенные шары $B_k$ исчезающего радиуса.
			\end{itemize}
			Пусть $x_k$ — центры этих шаров. По условию (3) выделим предел $x_k \to x_0$.
			Точка $x_0$ лежит в каком-то открытом $G_{\alpha_0}$. Значит, вокруг $x_0$ есть свободное место (шарик), целиком входящий в $G_{\alpha_0}$.
			
			\textbf{Финал:} С ростом $k$ наши «плохие» шары $B_k$ сжимаются к точке $x_0$ и неизбежно «проваливаются» в это свободное место. Значит, шар $B_k$ накрыт всего **одним** множеством $G_{\alpha_0}$. Противоречие с тем, что он «плохой».
		\end{enumerate}
		
		\item[\textbf{(3) $\Rightarrow$ (4)}] Пусть $M$ бесконечно. Выберем различные $x_n \in M$. По условию выделим сходящуюся подпоследовательность $x_{n_k} \to x_0$. Точка $x_0$ — предельная для $M$.
		
		\item[\textbf{(4) $\Rightarrow$ (3)}] Пусть $\{x_n\}$ — последовательность.
		\begin{itemize}
			\item Если множество значений конечно: выберем постоянную (сходящуюся) подпоследовательность.
			\item Если бесконечно: оно имеет предельную точку $x_0$. Выберем подпоследовательность, сходящуюся к $x_0$.
		\end{itemize}
	\end{itemize}
\end{proof}

\begin{remark}
	Условию компактности пространства $X$ также эквивалентно такое условие: любая непрерывная функция $f : X \to \mathbb{R}$ является ограниченной на $X$, то есть выполнено включение $C(X, \mathbb{R}) \subset B(X, \mathbb{R})$.
\end{remark}

\begin{theorem}[Кантора]\label{thm3.3}
	Пусть $X$ — компактное МП, и функция $f: X \to \mathbb{R}$ непрерывна на $X$. Тогда $f$ равномерно непрерывна на $X$.
\end{theorem}

\begin{definition}[Предкомпакт]
	Пусть $X$ — метрическое пространство. Множество $M \subset X$ называется \textbf{предкомпактом}, если множество $\overline M$ компактно как подпространство в $X$.
\end{definition}
\newpage % Был 6-й, стал 5-й
	\hypertarget{bilet6}{}\setupnav{5}{7} \subsection{Арцела-Асколи}

% --- Содержимое файла materials/3_3.tex ---

\begin{definition}[Пространство непрерывных функций]
	Обозначим за $C(X, Y)$ множество непрерывных функций $f: X \to Y$.
\end{definition}

\begin{theorem}[Арцела---Асколи]\label{arzela-ascoli}
	\textbf{Условия:}
	\begin{itemize}
		\item $X$ --- компактное метрическое пространство;
		\item $M \subset C(X, \mathbb{R})$.
	\end{itemize}
	\textbf{Следовательно:}
	\begin{itemize}
		\item Множество $M$ вполне ограниченно $\Leftrightarrow$ множество $M$ ограничено и выполнено условие \textbf{равностепенной непрерывности}:
		\[
		\forall \varepsilon > 0: \exists \delta > 0: \forall x, y \in X, \rho(x, y) < \delta: \forall f \in M: |f(x) - f(y)| < \varepsilon 
		\]
	\end{itemize}
\end{theorem}

\begin{idea}
	\textbf{Суть:} Для компактности в пространстве функций обычной ограниченности (все функции лежат в «трубе») мало. Нужно, чтобы функции не «осциллировали» слишком быстро — то есть для всех функций из набора должен существовать единый $\delta$, обеспечивающий малый шаг по значению.
\end{idea}

\begin{proof}[Доказательство для случая $X = [a, b]$]
	\begin{itemize}
		\item[$\Rightarrow$] \textbf{Необходимость.}
		\begin{enumerate}
			\item Ограниченность $M$ следует из существования конечной $\varepsilon$-сети (любое вполне ограниченное множество ограничено).
			\item Проверим равностепенную непрерывность. Зафиксируем $\varepsilon > 0$ и выберем конечную $\varepsilon$-сеть $\{\varphi_1, \dotsc, \varphi_n\} \subset C(X, \mathbb{R})$.
			\item Каждая $\varphi_k$ равномерно непрерывна (по теореме Кантора \ref{thm3.3}). Выберем $\delta_k > 0$ такие, что:
			\[ \rho(x, y) < \delta_k \Rightarrow |\varphi_k(x) - \varphi_k(y)| < \varepsilon. \]
			\item Положим $\delta := \min \{\delta_1, \dotsc, \delta_n\}$. Тогда для любой $f \in M$ выберем ближайшую $\varphi_k$ из сети и получим:
			\[ |f(x)-f(y)| \leqslant |f(x) - \varphi_k (x)| + |\varphi_k (x) - \varphi_k (y)| + |\varphi_k (y) - f (y)| < \varepsilon + \varepsilon + \varepsilon = 3\varepsilon. \]
		\end{enumerate}
		
		\item[$\Leftarrow$] \textbf{Достаточность.}
		\begin{enumerate}
			\item \textbf{Сетка значений:} Так как $M$ ограничено, $\exists C > 0 : \forall f \in M \Rightarrow \|f\| \leqslant C$. Зафиксируем $\varepsilon > 0$. По равностепенной непрерывности выберем $\delta > 0$.
			\item \textbf{Разбиение:} Разобьем $[a, b]$ точками $x_i$ на отрезки $< \delta$, а $[-C, C]$ точками $y_j$ на отрезки $< \varepsilon$. 
			\item \textbf{Построение сети:} Рассмотрим конечное множество $L$ кусочно-линейных функций (ломаных), проходящих через узлы $(x_i, y_j)$.
			
			
			
			\item \textbf{Аппроксимация:} Для любой $f \in M$ выберем ломаную $\varphi \in L$ такую, что в узлах $|f(x_i) - \varphi(x_i)| < \varepsilon$.
			\item \textbf{Оценка:} Пусть $x \in [x_i, x_{i+1}]$. Тогда:
			\[ |f(x) - \varphi(x)| \leqslant |f(x) - f(x_i)| + |f(x_i) - \varphi(x_i)| + |\varphi(x_i) - \varphi(x)| < 2\varepsilon + |\varphi(x_i) - \varphi(x_{i+1})|. \]
			Оценим шаг ломаной через значения $f$:
			\[ |\varphi(x_i) - \varphi(x_{i+1})| \leqslant |\varphi(x_i) - f(x_i)| + |f(x_i) - f(x_{i+1})| + |f(x_{i+1}) - \varphi(x_{i+1})| < 3\varepsilon. \]
			\item \textbf{Итог:} Таким образом, $\|f - \varphi\| < 5\varepsilon$. Множество $L$ конечно и образует $5\varepsilon$-сеть. Значит, $M$ вполне ограничено.
		\end{enumerate}
	\end{itemize}
\end{proof}

\begin{remark}
	Данная теорема является ключевым инструментом для доказательства существования решений дифференциальных уравнений (теорема Пеано).
\end{remark}\newpage % Был 7-й, стал 6-й
	
	% IV. ЛИНЕЙНЫЕ НОРМИРОВАННЫЕ ПРОСТРАНСТВА
	\section{Линейные нормированные пространства}
	\hypertarget{bilet7}{}\setupnav{6}{8} % --- Содержимое файла materials/4_1.tex ---

\label{seventh-bilet}
\subsection{Теорема Рисса}
\begin{definition}[Линейное нормированное пространство]
	Линейным нормированным пространством над полем $\mathbb{K}$ ($\mathbb{R}$ или $\mathbb{C}$) называется линейное пространство $E$ с функцией $\|\cdot\| : E \to \mathbb{R}$, обладающей свойствами:
	\begin{enumerate}
		\item $\forall x \in E: \|x\| \geqslant 0$, причем $\|x\| = 0 \Leftrightarrow x = 0$
		\item $\forall x \in E: \forall \alpha \in \mathbb{K}: \|\alpha x\| = |\alpha|\|x\|$
		\item $\forall x, y \in E: \|x + y\| \leqslant \|x\| + \|y\|$ (неравенство треугольника)
	\end{enumerate}
	Функция $\|\cdot\|$ называется \textbf{нормой} на пространстве $E$.
\end{definition}

\begin{example}
	Любое линейное нормированное пространство $E$ является метрическим с метрикой $\rho(x, y) := \|x - y\|$. Приведенные ранее метрические пространства являются линейными нормированными, и метрика в них имеет вид нормы разности элементов.
\end{example}

\begin{remark}
	Норма непрерывна как функция $E \to \mathbb{R}$. Если $x_n \xrightarrow[E]{} x$, то $\rho(x_n, x) = \|x_n - x\| \to 0$. По неравенству треугольника:
	\begin{itemize}
		\item $\|x_n\| \leqslant \|x_n - x\| + \|x\|$
		\item $\|x\| \leqslant \|x_n - x\| + \|x_n\|$
	\end{itemize}
	Таким образом, $\left|\|x_n\| - \|x\|\right| \to 0$, то есть $\|x_n\| \to \|x\|$.
\end{remark}

\begin{definition}
	Пусть $E$ --- линейное нормированное пространство. Множество $L \subset E$ называется:
	\begin{itemize}
		\item \textbf{Линейным многообразием}, если $L$ замкнуто относительно сложения и умножения на скаляры.
		\item \textbf{Подпространством}, если $L$ является линейным многообразием и при этом замкнуто как множество.
	\end{itemize}
\end{definition}

\begin{definition}[Линейная оболочка]
	Линейной оболочкой множества $M \subset E$ называется множество вида:
	\[[M] := \left\{\sum_{k = 1}^n \alpha_i m_i : \alpha_i \in \mathbb{K}, m_i \in M\right\}\]
\end{definition}

\begin{remark}
	Линейная оболочка всегда является линейным многообразием, но не всегда является подпространством.
\end{remark}

\begin{example}
	Пусть $E = C[0,1]$ и $M$ — многочлены. Тогда $M = [M]$ — многообразие, но не подпространство. По теореме Вейерштрасса, предел последовательности многочленов может не быть многочленом, значит $[M]$ не содержит все свои точки прикосновения (не замкнуто).
\end{example}

\begin{definition}[Эквивалентность норм]
	Нормы $\|\cdot\|$ и $\|\cdot\|'$ на $E$ называются \textbf{эквивалентными}, если $\exists C_1, C_2 > 0$ такие, что для любого $x \in E$:
	\[C_1\|x\|' \leqslant \|x\| \leqslant C_2\|x\|'\]
\end{definition}

\begin{definition}
	\textbf{Размерностью} пространства $E$ ($\dim E$) называется наибольший размер линейно независимой системы в $E$.
\end{definition}

\begin{proposition}[Лемма о «почти перпендикуляре»]\label{almostperp}
	Пусть $E$ --- ЛНП, $L \subsetneq E$ --- подпространство. Тогда для любого $\varepsilon > 0$ существует $y \in E$ такой, что $\|y\| = 1$ и $\rho(y, L) = \inf_{z \in L}\|y - z\| \geqslant 1 - \varepsilon$.
\end{proposition}

\begin{proof}
	\begin{enumerate}
		\item Зафиксируем $y_0 \in E \setminus L$ и положим $d := \rho(y_0, L) > 0$.
		\item Выберем $z_0 \in L$ такой, что $d \leqslant \|y_0 - z_0\| \leqslant \frac{d}{1-\varepsilon}$. 
		\item Покажем, что подходит вектор $y := \frac{y_0 - z_0}{\|y_0 - z_0\|}$. Очевидно, $\|y\| = 1$.
		\item Положим $\alpha = \frac{1}{\|y_0 - z_0\|}$. Для любого $z \in L$:
		\[ \|y - z\| = \|\alpha(y_0 - z_0) - z\| = \alpha \| y_0 - (z_0 + \frac{1}{\alpha} z) \| \]
		\item Так как $(z_0 + \frac{1}{\alpha} z) \in L$, то $\|y_0 - (z_0 + \frac{1}{\alpha} z)\| \geqslant d$. Тогда:
		\[ \|y - z\| \geqslant \alpha d = \frac{d}{\|y_0 - z_0\|} \geqslant \frac{d}{d/(1-\varepsilon)} = 1 - \varepsilon \]
		Значит, $\inf_{z \in L}\|y - z\| \geqslant 1 - \varepsilon$.
	\end{enumerate}
\end{proof}

\begin{idea}[Лемма о почти перпендикуляре]
	\textbf{Геометрический смысл:} В конечномерном случае мы всегда можем найти вектор, строго перпендикулярный плоскости. В произвольном ЛНП «идеального» перпендикуляра может не существовать, но лемма утверждает, что мы всегда можем найти вектор $y$, который «почти» перпендикулярен подпространству $L$.
	
	\textbf{Важное уточнение по расстоянию:}
	\begin{itemize}
		\item Мы измеряем расстояние от \textbf{конца (носка)} вектора $y$ до подпространства $L$. 
		\item Хотя «начало» вектора (пятка) находится в нуле и принадлежит $L$, сам вектор «торчит» в сторону сферы. 
		\item Так как длина вектора равна $1$, кратчайший путь от его конца до $L$ \textbf{не может быть больше 1} (путь в ноль всегда доступен).
		\item Лемма гарантирует, что мы можем направить вектор так, чтобы его конец был удален от $L$ \textbf{почти на всю свою длину} ($\geqslant 1 - \varepsilon$).
	\end{itemize}
	
	\textbf{Суть доказательства:} Мы берем вектор $y_0$ вне подпространства и «нормируем» его отклонение от почти ближайшей точки $z_0 \in L$. Чем ближе $z_0$ к истинному инфимуму, тем «вертикальнее» стоит итоговый вектор $y$ относительно $L$.
\end{idea}

\begin{theorem}[Теорема Рисса]\label{thm4.1}
	Пусть $E$ --- ЛНП. Тогда единичная сфера $S(0, 1)$ компактна в $E \Leftrightarrow \dim E < +\infty$.
\end{theorem}

\begin{proof}
	\begin{itemize}
		\item[$\Leftarrow$] В конечномерном пространстве все нормы эквивалентны. Относительно евклидовой нормы сфера $S(0, 1)$ компактна по теореме Гейне-Бореля.
		
		\item[$\Rightarrow$] \textbf{От противного.} Пусть $\dim E = \infty$. Построим последовательность $\{x_n\} \subset S(0, 1)$ с расстояниями $\rho(x_i, x_j) \geqslant 1 - \varepsilon$:
		\begin{enumerate}
			\item Выберем $x_1 \in S(0, 1)$ произвольно.
			\item Пусть $L_1 = [x_1]$. Т.к. $\dim L_1 < \infty$, оно замкнуто и является подпространством. По лемме~\ref{almostperp} выберем $x_2 \in S(0, 1)$ так, что $\rho(x_2, L_1) \geqslant 1 - \varepsilon$.
			\item Пусть $L_2 = [x_1, x_2]$. Оно также замкнуто. Выберем $x_3 \in S(0, 1)$ так, что $\rho(x_3, L_2) \geqslant 1 - \varepsilon$.
			\item Процесс не закончится, так как $\dim E = \infty$. Получим последовательность, из которой нельзя выделить сходящуюся (т.к. $\rho(x_n, x_m) \geqslant 1-\varepsilon$). Значит, сфера не вполне ограничена и не компактна. Противоречие.
		\end{enumerate}
	\end{itemize}
\end{proof}	

\begin{idea}[Смысл теоремы Рисса]
	\textbf{Суть:} Это критерий того, насколько «просторно» в пространстве. 
	
	\begin{itemize}
		\item \textbf{В конечномерном мире} (как $\mathbb{R}^3$) сфера «тесная». Если наставить на ней бесконечно много точек, им придется прижиматься друг к другу (сгущаться). Это обеспечивает компактность.
		\item \textbf{В бесконечномерном мире} всегда есть «куда расти». Там можно расставить бесконечный «ёж» из векторов так, что все они будут на расстоянии 1 друг от друга и никогда не сблизятся. 
	\end{itemize}
	\textbf{Вывод:} Сфера компактна только тогда, когда пространство конечно. Бесконечномерность «убивает» компактность из-за бесконечного избытка места.
\end{idea}

\begin{idea}[Теорема Рисса]
	\textbf{Почему это важно:} Эта теорема проводит жесткую границу между привычным нам миром (конечномерным) и миром бесконечных функций. 
	
	\textbf{Геометрия бесконечномерности:}
	\begin{itemize}
		\item В конечномерном пространстве «место» ограничено. Если мы пытаемся расставить бесконечно много точек на сфере, им приходится сгущаться, образуя предельную точку.
		\item В бесконечномерном пространстве появляется «бесконечно много новых направлений». Благодаря лемме о почти перпендикуляре, мы можем каждый раз выбирать вектор, который \textbf{почти перпендикулярен всем предыдущим}.
		\item В итоге мы строим последовательность $\{x_n\}$, где расстояние между любыми двумя точками $\rho(x_n, x_m) \approx 1$. 
	\end{itemize}
	
	\textbf{Суть противоречия:} 
	Такая последовательность «разлетается» во все стороны. Она никогда не сможет сблизиться (стать фундаментальной), а значит, из нее нельзя выделить сходящуюся подпоследовательность. Сфера оказывается «слишком просторной», чтобы быть компактной.
\end{idea}\newpage
	\hypertarget{bilet8}{}\setupnav{7}{9} % --- Содержимое файла materials/5_1.tex ---
\subsection{Характеристическое свойство евклидовых пространств. Ба-
	наховы и гильбертовы пространства}
\begin{definition}[Евклидово пространство]
	\textbf{Евклидовым пространством} над полем $\mathbb{K}$ (где $\mathbb{K} = \mathbb{R}$ или $\mathbb{K} = \mathbb{C}$) называется линейное пространство $E$ над $\mathbb{K}$ с функцией $(\cdot, \cdot) : E^2 \to \mathbb{K}$, называемой \textbf{скалярным произведением}, которая обладает следующими свойствами:
	\begin{enumerate}
		\item $\forall x \in E: (x, x) \geqslant 0$, причем $(x, x) = 0 \Leftrightarrow x = 0$ (положительная определенность);
		\item $\forall x, y \in E: (x, y) = \overline{(y, x)}$ (эрмитова симметричность);
		\item $\forall x, y, z \in E, \forall \alpha, \beta \in \mathbb{K}: (\alpha x + \beta y, z) = \alpha(x, z) + \beta(y, z)$ (линейность по первому аргументу).
	\end{enumerate}
\end{definition}

\begin{theorem}[Характеристическое свойство Евклидовых пространств]
	Пусть $E$ --- линейное нормированное пространство. Норма в $E$ порождается скалярным произведением тогда и только тогда, когда для любых $x, y \in E$ выполняется \textbf{равенство параллелограмма}:
	\[\|x+y\|^2 + \|x-y\|^2 = 2\|x\|^2 + 2\|y\|^2\]
\end{theorem}


\begin{definition}[Типы пространств]~
	\begin{itemize}
		\item Полное линейное нормированное пространство $E$ называется \textbf{банаховым}.
		\item Полное евклидово пространство $H$ называется \textbf{гильбертовым}.
	\end{itemize}
\end{definition}

\begin{corollary}[Неравенство Коши---Буняковского]
	Пусть $E$ --- евклидово пространство. Тогда для любых $x, y \in E$:
	\[|(x, y)| \leqslant \|x\| \cdot \|y\|\]
\end{corollary}

\begin{corollary}
	Из неравенства Коши-Буняковского следует непрерывность скалярного произведения. Если $x_n \to x_0$ и $y_n \to y_0$, то:
	\[ |(x_n, y_n) - (x_0, y_0)| \leqslant |(x_n - x_0, y_n)| + |(x_0, y_n - y_0)| \leqslant \|x_n - x_0\|\|y_n\| + \|x_0\|\|y_n - y_0\| \to 0 \]
\end{corollary}\newpage
	\hypertarget{bilet9}{}\setupnav{8}{10} % --- Содержимое файла materials/5_2.tex ---
\subsection{Эквивалентность норм в конечномерном пространстве. Понятие линейного топологического пространства, примеры}
\begin{definition}[Линейное пространство]
	Непустое множество $L$ называется \textbf{линейным пространством} над полем $\mathbb{K}$, если оно удовлетворяет следующим условиям:
	\begin{enumerate}
		\item Для любых $x, y \in L$ однозначно определена сумма $x + y \in L$:
		\begin{enumerate}
			\item $x + y = y + x$ (коммутативность);
			\item $x + (y + z) = (x + y) + z$ (ассоциативность);
			\item $\exists 0 \in L : \forall x \in L \Rightarrow x + 0 = x$ (существование нуля);
			\item $\forall x \in L, \exists (-x) \in L : x + (-x) = 0$ (существование обратного элемента).
		\end{enumerate}
		\item Для любого $\alpha \in \mathbb{K}$ и $x \in L$ определено произведение $\alpha x \in L$:
		\begin{enumerate}
			\item $\alpha (\beta x) = (\alpha \beta) x$ (ассоциативность умножения);
			\item $1 \cdot x = x$ (унитарность);
			\item $(\alpha + \beta) x = \alpha x + \beta x$ (дистрибутивность по скалярам);
			\item $\alpha (x + y) = \alpha x + \alpha y$ (дистрибутивность по векторам).
		\end{enumerate}
	\end{enumerate}
\end{definition}

\begin{definition}[Изоморфизмы]
	\begin{itemize}
		\item Линейные пространства $L$ и $L^*$ называются \textbf{изоморфными}, если существует биекция $f: L \to L^*$, сохраняющая линейные операции:
		\[ x \leftrightarrow x^*, y \leftrightarrow y^* \implies x+y \leftrightarrow x^*+y^*, \quad \alpha x \leftrightarrow \alpha x^*. \]
		\item Евклидовы пространства $E$ и $E^*$ называются \textbf{изоморфными}, если они изоморфны как линейные и сохраняют скалярное произведение:
		\[ x \leftrightarrow x^*, y \leftrightarrow y^* \implies (x, y)_E = (x^*, y^*)_{E^*}. \]
	\end{itemize}
\end{definition}

\begin{remark}
	Смежные определения (нормированного пространства, скалярного произведения) сформулированы в билете №7.
\end{remark}

\begin{example}[Эквивалентные нормы в пространстве функций]
	Рассмотрим $X = \{u \in C^1[0, 1] : u(0) = 0\}$. Введем две нормы:
	\[ \|u\|_1 = \left(\int_0^1 (u^2 + (u')^2) dx\right)^{1/2}, \quad \|u\|_2 = \left(\int_0^1 (u')^2 dx\right)^{1/2} \]
	Докажем их эквивалентность ($C_1 \|u\|_2 \leqslant \|u\|_1 \leqslant C_2 \|u\|_2$).
	\begin{enumerate}
		\item Очевидно, что $\|u\|_2 \leqslant \|u\|_1$ (так как подынтегральное выражение первой нормы больше).
		\item Для обратного неравенства используем $u(0)=0$: $u(x) = \int_0^x u'(t) dt$.
		\item Возведем в квадрат и применим неравенство Коши-Буняковского:
		\[ u^2(x) = \left(\int_0^x 1 \cdot u'(t)dt \right)^2 \leqslant \left(\int_0^x 1^2 dt\right) \left(\int_0^x (u'(t))^2 dt\right) \leqslant x \cdot \|u\|_2^2 \leqslant \|u\|_2^2. \]
		\item Интегрируем по $x$ от 0 до 1:
		\[ \int_0^1 u^2(x) dx \leqslant \int_0^1 \|u\|_2^2 dx = \|u\|_2^2. \]
		\item Итого: $\|u\|_1^2 = \int u^2 + \int (u')^2 \leqslant \|u\|_2^2 + \|u\|_2^2 = 2\|u\|_2^2$.
		Отсюда $\|u\|_1 \leqslant \sqrt{2}\|u\|_2$. Нормы эквивалентны.
	\end{enumerate}
\end{example}



\begin{example}[Неэквивалентные нормы]
	Рассмотрим $C[a, b]$ с нормами $\|x\|_C = \max |x(t)|$ и $\|x\|_1 = \int_a^b |x(t)| dt$.
	\begin{itemize}
		\item Неравенство $\|x\|_1 \leqslant (b-a)\|x\|_C$ выполнено всегда.
		\item Обратное неравенство невозможно. Построим последовательность «пиков»:
		\[
		x_n(t) = \begin{cases} 
			2n^2(t-a), & t \in [a, a + \frac{1}{2n}] \\
			2n - 2n^2(t - (a+\frac{1}{2n})), & t \in [a+\frac{1}{2n}, a+\frac{1}{n}] \\
			0, & \text{иначе}
		\end{cases}
		\]
		Площадь под графиком (интеграл) $\|x_n\|_1 = \frac{1}{2} \cdot \frac{1}{n} \cdot 2n = 1$ (константа).
		Максимум модуля $\|x_n\|_C = 2n \to \infty$.
		Константы $K$, такой что $2n \leqslant K \cdot 1$ для всех $n$, не существует.
	\end{itemize}
\end{example}



\begin{proposition}[Эквивалентность норм в конечномерном пространстве]\label{normequiv}
	Пусть $E$ --- ЛНП и $\dim E < +\infty$. Тогда любые две нормы на $E$ эквивалентны.
\end{proposition}



\begin{proof}
	Выберем базис $e = \{e_1, \dots, e_n\}$. Любой $x$ задается столбцом координат $\alpha$.
	Введем \textit{евклидову норму}: $\|x\|_{euc} := \sqrt{\sum \alpha_k^2}$.
	Докажем эквивалентность произвольной нормы $\|\cdot\|$ и $\|\cdot\|_{euc}$.
	
	\begin{enumerate}
		\item \textbf{Оценка сверху} ($\|\cdot\| \leqslant C \|\cdot\|_{euc}$).
		\[ \|x\| = \left\| \sum \alpha_k e_k \right\| \leqslant \sum |\alpha_k| \|e_k\| \leqslant \left(\max \|e_k\|\right) \sum |\alpha_k|. \]
		Используя неравенство Коши для сумм ($\sum |\alpha_k| \leqslant \sqrt{n} \sqrt{\sum \alpha_k^2}$), получаем $C = \sqrt{n} \max \|e_k\|$.
		
		\item \textbf{Оценка снизу} ($\|\cdot\|_{euc} \leqslant \widetilde{C} \|\cdot\|$). Докажем от противного.
		\begin{itemize}
			\item Пусть константы не существует. Тогда $\forall n \in \mathbb{N} \exists x_n$: $\|x_n\|_{euc} > n \|x_n\|$.
			\item Нормируем векторы: пусть $y_n = x_n / \|x_n\|_{euc}$. Тогда $\|y_n\|_{euc} = 1$, а $\|y_n\| < 1/n$.
			\item Единичная сфера в евклидовом пространстве компактна. Выделим сходящуюся подпоследовательность: $y_{n_k} \xrightarrow{\|\cdot\|_{euc}} y_0$, где $\|y_0\|_{euc} = 1$.
			\item В силу пункта (1), сходимость по евклидовой норме влечет сходимость по норме $\|\cdot\|$:
			\[ \|y_{n_k} - y_0\| \leqslant C \|y_{n_k} - y_0\|_{euc} \to 0. \]
			\item Значит, $\|y_{n_k}\| \to \|y_0\|$. Но по построению $\|y_{n_k}\| < 1/n_k \to 0$.
			\item Получаем $\|y_0\| = 0 \Rightarrow y_0 = 0$.
			\item \textbf{Противоречие:} $\|y_0\|_{euc} = 1 \neq 0$.
		\end{itemize}
	\end{enumerate}
\end{proof}

\begin{corollary}\label{closeness-subspace}
	Любое конечномерное подпространство $L \subset E$ замкнуто.
\end{corollary}

\begin{proof}
	Так как $\dim L < \infty$, все нормы на нем эквивалентны. Относительно евклидовой нормы $L$ изоморфно $\mathbb{R}^n$, а $\mathbb{R}^n$ полно.
	Полное подпространство в метрическом пространстве всегда замкнуто.
\end{proof}

\begin{definition}[Линейное топологическое пространство]
	Топологическое пространство $X$ над полем $\mathbb{K}$, в котором операции сложения и умножения на скаляр непрерывны.
\end{definition}\newpage
	\hypertarget{bilet10}{}\setupnav{9}{11} % --- Содержимое файла materials/4_4.tex ---

\subsection{Элемент наилучшего приближения. Теорема о проекции}

\begin{definition}[Элемент наилучшего приближения]
	Пусть $E$ --- линейное нормированное пространство, $L \subset E$ --- подпространство, $h \in E$. \textbf{Элементом наилучшего приближения} для $h$ называется вектор $x \in L$ такой, что расстояние от $h$ до $x$ минимально среди всех векторов $L$:
	\[ \|h - x\| = \rho(h, L) = \inf \limits_{y \in L} \|h - y\|. \]
\end{definition}

\begin{remark}
	В общем случае элемент наилучшего приближения может не существовать (если подпространство не замкнуто или пространство не рефлексивно) или быть не единственным (если норма "угловатая").
\end{remark}

\begin{proposition}[Существование и единственность в Гильбертовом пространстве]\label{bestelt}
	Пусть $H$ --- гильбертово пространство, $M \subset H$ --- \textbf{замкнутое} подпространство. Тогда для любого $h \in H$ существует единственный элемент наилучшего приближения $x \in M$.
\end{proposition}

\begin{proof}
	Зафиксируем $h \in H$. Обозначим $d := \rho(h, M)$.
	
	\textbf{1. Существование.}
	\begin{itemize}
		\item Если $d = 0$, то в силу замкнутости $M$, $h \in M$, и искомый элемент — сам $h$.
		\item Пусть $d > 0$. По определению инфимума выберем минимизирующую последовательность $\{x_n\} \subset M$ такую, что $\|h - x_n\| \to d$.
		\item Проверим фундаментальность $\{x_n\}$. Применим равенство параллелограмма к векторам $(h - x_n)$ и $(h - x_m)$:
		\[ \|(h - x_n) + (h - x_m)\|^2 + \|(h - x_n) - (h - x_m)\|^2 = 2\|h - x_n\|^2 + 2\|h - x_m\|^2 \]
		Преобразуем:
		\[ \|2h - (x_n + x_m)\|^2 + \|x_m - x_n\|^2 = 2\|h - x_n\|^2 + 2\|h - x_m\|^2 \]
		\[ 4\left\|h - \frac{x_n + x_m}{2}\right\|^2 + \|x_n - x_m\|^2 = 2\|h - x_n\|^2 + 2\|h - x_m\|^2 \]
		\item Заметим, что $\frac{x_n + x_m}{2} \in M$, следовательно $\left\|h - \frac{x_n + x_m}{2}\right\| \geqslant d$.
		\item Переходим к пределу. Правая часть стремится к $2d^2 + 2d^2 = 4d^2$. Левая часть $\geqslant 4d^2 + \|x_n - x_m\|^2$. Отсюда следует, что $\|x_n - x_m\| \to 0$.
		\item Так как $H$ полно и $M$ замкнуто, $x_n \to x \in M$. В силу непрерывности нормы $\|h - x\| = d$.
	\end{itemize}
	
	\textbf{2. Единственность.}
	\begin{itemize}
		\item Пусть существуют два элемента $x, y \in M$, реализующих минимум: $\|h - x\| = d$ и $\|h - y\| = d$.
		\item Применим равенство параллелограмма к векторам $h-x$ и $h-y$:
		\[ \|(h-x) + (h-y)\|^2 + \|(h-x) - (h-y)\|^2 = 2\|h-x\|^2 + 2\|h-y\|^2 \]
		\[ 4\left\|h - \frac{x+y}{2}\right\|^2 + \|y-x\|^2 = 2d^2 + 2d^2 = 4d^2 \]
		\item Так как $\frac{x+y}{2} \in M$, то $\left\|h - \frac{x+y}{2}\right\| \geqslant d$, а значит первое слагаемое $\geqslant 4d^2$.
		\item Получаем неравенство $4d^2 + \|x-y\|^2 \leqslant 4d^2$, откуда $\|x - y\| = 0 \Rightarrow x = y$.
	\end{itemize}
\end{proof}

\begin{idea}[О наилучшем приближении]
	
	\textbf{Смысл утверждения:}
	
	Представьте, что $M$ — это бесконечная плоская поверхность (подпространство), а $h$ — точка, висящая над ней. Утверждение гарантирует, что на этой поверхности есть \textbf{ровно одна} точка $x$, ближайшая к $h$ (проекция).
	
	\begin{itemize}
		
		\item В «плохих» пространствах (как $L_\infty$ с квадратными шарами) ближайших точек может быть много (целая грань квадрата).
		
		\item В Гильбертовом пространстве шары «круглые» (строго выпуклые), поэтому касание происходит только в одной точке.
		
	\end{itemize}
	
	\textbf{Идея доказательства:}
	
	\begin{enumerate}
		
		\item \textbf{Существование:} Мы берем последовательность точек, которые подходят всё ближе к идеальному расстоянию $d$. Чтобы доказать, что они не «танцуют» на месте, а сходятся к цели, мы используем \textbf{равенство параллелограмма}. Оно переводит геометрическое сближение (расстояния) в алгебраический факт фундаментальности.
		
		\item \textbf{Единственность:} Если бы было две наилучшие точки $x$ и $y$, то середина отрезка между ними $\frac{x+y}{2}$ оказалась бы \textit{еще ближе} к цели (из-за выпуклости шара), что противоречит тому, что $x$ и $y$ — уже самые близкие.
		
	\end{enumerate}
	
\end{idea}

\begin{definition}[Аннулятор]
	Пусть $E$ --- евклидово пространство, $S \subset E$. \textbf{Аннулятором} (или ортогональным дополнением) множества $S$ называется множество:
	\[S^\perp := \{y \in E \mid \forall x \in S: (x, y) = 0\}\]
\end{definition}

\begin{remark}
	$S^\perp$ всегда является \textbf{замкнутым подпространством} в $E$ (даже если $S$ не замкнуто и не подпространство).
	Справедливо: $S^\perp = [S]^\perp = \overline{[S]}^\perp$.
\end{remark}

\begin{theorem}[Теорема Рисса о проекции]\label{thm4.3}
	Пусть $H$ --- гильбертово пространство, $M \subset H$ --- \textbf{замкнутое} подпространство. Тогда пространство разлагается в прямую ортогональную сумму:
	\[ H = M \oplus M^\perp \]
	Это означает, что любой вектор $h \in H$ единственным образом представим в виде $h = x + y$, где $x \in M, y \in M^\perp$.
\end{theorem}

\begin{idea}
	Теорема утверждает, что мы можем "уронить перпендикуляр" из любой точки на подпространство. Вектор $x$ — это проекция ("тень"), а вектор $y$ — перпендикуляр ("высота").
\end{idea}


\begin{proof}
	\textbf{1. Возможность разложения ($H = M + M^\perp$).}
	\begin{itemize}
		\item Зафиксируем $h \in H$. По предложению \ref{bestelt}, существует единственный элемент наилучшего приближения $x \in M$.
		\item Обозначим вектор невязки $y := h - x$. Покажем, что $y \in M^\perp$.
		\item Обозначим $d = \|y\| = \|h - x\|$. Для любого вектора $m \in M$ и любого числа $\alpha \in \mathbb{R}$ вектор $(x + \alpha m)$ лежит в $M$.
		\item Так как $x$ — ближайший элемент, расстояние до него минимально:
		\[ \|y\|^2 \leqslant \|h - (x + \alpha m)\|^2 = \|y - \alpha m\|^2 = (y - \alpha m, y - \alpha m) \]
		\[ \|y\|^2 \leqslant \|y\|^2 - 2\alpha(y, m) + \alpha^2\|m\|^2 \]
		\item Сокращаем $\|y\|^2$ и получаем: $2\alpha(y, m) \leqslant \alpha^2\|m\|^2$.
		\item Если бы $(y, m) \neq 0$, мы могли бы выбрать $\alpha$ очень малым и такого знака, что левая часть (линейная по $\alpha$) перевесила бы правую (квадратичную по $\alpha$), нарушив неравенство. (Формально: положим $\alpha = \text{sign}(y,m) \cdot \varepsilon$, при $\varepsilon \to 0$ получаем противоречие).
		\item Следовательно, $(y, m) = 0$ для всех $m \in M$. Значит $y \perp M$.
		\item Мы представили $h = x + y$, где $x \in M, y \in M^\perp$.
	\end{itemize}
	
	\textbf{2. Единственность (Прямая сумма).}
	\begin{itemize}
		\item Пусть $z \in M \cap M^\perp$.
		\item Так как $z \in M^\perp$ и $z \in M$, он должен быть ортогонален самому себе: $(z, z) = 0$.
		\item По аксиоме скалярного произведения $\|z\|^2 = 0 \Rightarrow z = 0$.
		\item Пересечение тривиально, значит сумма прямая.
	\end{itemize}
\end{proof}

\begin{remark}
	Это свойство (существование проекции) выделяет гильбертовы пространства среди всех банаховых. В общем случае в банаховом пространстве нельзя гарантировать существование такого "удобного" разложения.
\end{remark}

\begin{idea}[О теореме Рисса]
	
	\textbf{Смысл утверждения:}
	
	Любой вектор можно разложить на две независимые составляющие: «тень» на подпространстве ($x \in M$) и «высоту» над ним ($y \in M^\perp$). Это обобщение школьного разложения сил на нормальную и тангенциальную составляющие.
	
	\textbf{Идея доказательства (Вариационный метод):}
	
	Мы уже знаем, что $x$ — это самая близкая точка в $M$ к вектору $h$. Вектор ошибки (высота) — это $y = h - x$.
	
	Почему $y$ обязан быть перпендикулярен $M$?
	
	\begin{itemize}
		
		\item Представьте, что вектор $y$ наклонен к плоскости $M$ под углом $80^\circ$ (не перпендикулярен).
		
		\item Тогда мы могли бы чуть-чуть сместиться по плоскости в сторону этого наклона и уменьшить длину вектора $y$.
		
		\item Но $x$ — это уже \textbf{самая} близкая точка, уменьшать расстояние некуда.
		
		\item Значит, никакого наклона нет, угол ровно $90^\circ$.
		
	\end{itemize}
	
\end{idea}\newpage
	\hypertarget{bilet11}{}\setupnav{10}{12} % --- Содержимое файла materials/4_5.tex ---

\subsection{Базисы в гильбертовых пространствах. Равенство Парсеваля}

\begin{definition}[Базис Шаудера]
	Пусть $E$ --- линейное нормированное пространство. Система $\{e_n\} \subset E$ называется \textbf{базисом Шаудера} в $E$, если для любого $x \in E$ существует \textit{единственный} набор чисел $\{\alpha_n\} \subset \mathbb{K}$ такой, что ряд сходится к $x$ по норме пространства:
	\[ x = \sum_{n = 1}^{\infty} \alpha_n e_n \]
\end{definition}

\begin{proposition}[Неравенство Бесселя]
	Пусть $E$ --- евклидово пространство, $\{e_i\}$ --- \textbf{ортонормированная система}. Тогда для любого $x \in E$:
	\[\sum_{k = 1}^{\infty}|(x, e_k)|^2 \leqslant \|x\|^2\]
\end{proposition}
\begin{proof}[Доказательство неравенства Бесселя]
	Достаточно заметить, что в силу ортонормированности системы $\{e_k\}$ выполнено равенство
	$$
	0 \le \left\|x-\sum_{k=1}^{n}(x,e_k)e_k\right\|^2
	= \|x\|^2-\sum_{k=1}^{n}|(x,e_k)|^2.
	$$
	
	Поскольку левая часть неотрицательна, то
	$$
	\sum_{k=1}^{n}|(x,e_k)|^2 \le \|x\|^2
	\quad \text{для любого } n\ge1.
	$$
	
	Переходя к пределу при $n\to\infty$, получаем
	$$
	\sum_{k=1}^{\infty}|(x,e_k)|^2 \le \|x\|^2.
	$$
\end{proof}

\begin{theorem}[Неравенство Коши--Буняковского]
	Пусть $E$ --- евклидово пространство. Тогда для любых $x,y\in E$ выполнено неравенство
	$$
	|(x,y)|\le \|x\|\,\|y\|.
	$$
\end{theorem}

\begin{proof}
	Случай $y=0$ тривиален, поэтому можно считать, что $y\ne 0$. Тогда система
	$$
	\left\{\frac{y}{\|y\|}\right\}
	$$
	является ортонормированной. По неравенству Бесселя для любого $x\in E$ имеем
	$$
	\left|\left(x,\frac{y}{\|y\|}\right)\right|^2 \le \|x\|^2.
	$$
	Но
	$$
	\left(x,\frac{y}{\|y\|}\right)=\frac{(x,y)}{\|y\|},
	$$
	следовательно,
	$$
	\frac{|(x,y)|^2}{\|y\|^2}\le \|x\|^2.
	$$
	Умножая на $\|y\|^2$ и извлекая квадратный корень, получаем
	$$
	|(x,y)|\le \|x\|\,\|y\|.
	$$
\end{proof}

\begin{proposition}[Экстремальное свойство коэффициентов Фурье]\label{minproperty}
	Пусть $\{e_k\}_{k=1}^n$ --- ортонормированная система в евклидовом пространстве $E$. Тогда для любого $x \in E$ и любых чисел $\alpha_k$:
	\[\left\|x - \sum_{k = 1}^n \alpha_ke_k\right\| \geqslant \left\|x - \sum_{k = 1}^n (x, e_k)e_k\right\|\]
	Равенство достигается тогда и только тогда, когда $\alpha_k = (x, e_k)$ (коэффициенты Фурье).
\end{proposition}

\begin{proof}
	[Доказательство для случая, когда $\mathbb{K}=\mathbb{R}$]
	В силу ортонормированности системы $\{e_k\}_{k=1}^n$ имеем
	$$
	\left\|x-\sum_{k=1}^n \alpha_k e_k\right\|^2
	=
	\left(
	x-\sum_{k=1}^n \alpha_k e_k,\;
	x-\sum_{k=1}^n \alpha_k e_k
	\right).
	$$
	
	Раскрывая скалярное произведение, получаем
	$$
	=
	\|x\|^2
	-2\sum_{k=1}^n \alpha_k (x,e_k)
	+\sum_{k=1}^n \alpha_k^2.
	$$
	
	С другой стороны,
	$$
	\left\|x-\sum_{k=1}^n (x,e_k)e_k\right\|^2
	=
	\|x\|^2-\sum_{k=1}^n (x,e_k)^2.
	$$
	
	Поэтому
	$$
	\left\|x-\sum_{k=1}^n \alpha_k e_k\right\|^2
	=
	\left\|x-\sum_{k=1}^n (x,e_k)e_k\right\|^2
	+\sum_{k=1}^n\big((x,e_k)-\alpha_k\big)^2.
	$$
	
	Отсюда следует неравенство
	$$
	\left\|x-\sum_{k=1}^n \alpha_k e_k\right\|
	\ge
	\left\|x-\sum_{k=1}^n (x,e_k)e_k\right\|.
	$$
	
	Причём равенство достигается тогда и только тогда, когда
	$$
	\alpha_k=(x,e_k),\quad k=1,\dots,n.
	$$
\end{proof}

\begin{idea}
	Это алгебраическое выражение геометрического факта: перпендикуляр короче наклонной. Проекция вектора $x$ на подпространство, натянутое на $e_k$, дается именно коэффициентами Фурье. Любой другой выбор коэффициентов $\alpha_k$ отклоняет нас от проекции.
\end{idea}

\begin{definition}[Сепарабельное пространство]\label{separability}
	Метрическое пространство $X$ называется \textbf{сепарабельным}, если в нем существует не более чем счетное всюду плотное множество.
	\[ \exists S \subset X : |S| \le \aleph_0 \quad \text{и} \quad \overline{S} = X \]
\end{definition}

\begin{idea}
	\textbf{Суть:} Сепарабельность означает, что пространство хоть и «большое» (несчетное), но его скелет — «маленький» (счетный). 
	Это позволяет нам контролировать пространство, используя последовательности (например, приближать любой элемент пределом последовательности из счетного набора).
	
	\textbf{Примеры:}
	\begin{itemize}
		\item $\mathbb{R}$ — сепарабельно (множество $\mathbb{Q}$ счетно и плотно).
		\item $C[a, b]$ — сепарабельно (многочлены с рациональными коэффициентами).
		\item $l_\infty$ — \textbf{не} сепарабельно (там слишком много «изолированных» точек).
	\end{itemize}
\end{idea}

\begin{theorem}[Критерий базиса]\label{thm4.4}
	Пусть $H$ --- сепарабельное гильбертово пространство, $\dim{H} = \infty$, и $e = \{e_n\}$ --- ортонормированная система. Следующие условия эквивалентны:
	\begin{enumerate}
		\item $e$ --- ортонормированный базис;
		\item Система $e$ полна (замкнутая линейная оболочка совпадает с $H$: $\overline{[e]} = H$);
		\item Для любого $h \in H$ выполнено \textbf{равенство Парсеваля}: $\|h\|^2 = \sum_{n = 1}^\infty |(h, e_n)|^2$;
		\item Аннулятор системы тривиален: $e^\perp = \{0\}$ (нет ненулевого вектора, ортогонального всем $e_n$).
	\end{enumerate}
\end{theorem}

\begin{proof}
	\begin{itemize}
		\item $(1) \Rightarrow (2)$
		По определению базиса, любой $x$ есть предел частичных сумм ряда $\sum \alpha_k e_k$, а частичные суммы лежат в $[e]$. Значит $x \in \overline{[e]}$.
		
		\item $(2) \Rightarrow (1)$
		\begin{enumerate}
			\item \textbf{Существование:} Зафиксируем $h \in H, \varepsilon > 0$. Так как $\overline{[e]} = H$, существует линейная комбинация $\sum_{k=1}^n \alpha_k e_k$, приближающая $h$ с точностью $\varepsilon$.
			\item По предложению \ref{minproperty}, сумма с коэффициентами Фурье приближает еще лучше:
			\[ \left\|h - \sum_{k=1}^n (h, e_k)e_k\right\| \leqslant \left\|h - \sum_{k=1}^n \alpha_k e_k\right\| < \varepsilon. \]
			\item Значит, ряд Фурье сходится к $h$.
			\item \textbf{Единственность:} Если $h = \sum \beta_n e_n$, то скалярно умножив на $e_k$ (и пользуясь непрерывностью скалярного произведения), получим $(h, e_k) = \sum \beta_n (e_n, e_k) = \beta_k$. Коэффициенты определены однозначно.
		\end{enumerate}
		
		\item $(1) \Leftrightarrow (3)$
		Ранее показано: $\left\|h - \sum_{k=1}^n (h, e_k)e_k\right\|^2 = \|h\|^2 - \sum_{k=1}^n |(h, e_k)|^2$.
		Сходимость ряда к $h$ (левая часть $\to 0$) равносильна сходимости суммы квадратов коэффициентов к квадрату нормы (правая часть $\to 0$).
		
		\item $(2) \Leftrightarrow (4)$
		По теореме о проекции, $H = \overline{[e]} \oplus \overline{[e]}^\perp$.
		Так как $e^\perp = [e]^\perp = \overline{[e]}^\perp$, то условие $\overline{[e]} = H$ равносильно $e^\perp = \{0\}$.
	\end{itemize}
\end{proof}

\begin{corollary}
	В любом бесконечномерном сепарабельном гильбертовом пространстве $H$ существует счетный ортонормированный базис.
\end{corollary}

\begin{idea}[Смысл теоремы]
	
	Эта теорема утверждает, что «хороший» базис в бесконечномерном пространстве ведет себя так же, как обычный базис $(\vec{i}, \vec{j}, \vec{k})$ в школьной геометрии.
	
	Все 4 условия говорят об одном и том же — \textbf{о полноте набора векторов}:
	
	\begin{itemize}
		
		\item \textbf{(1) Базис:} Любой вектор можно собрать из кирпичиков $e_n$.
		
		\item \textbf{(2) Полнота системы:} Линейными комбинациями $e_n$ можно подобраться сколь угодно близко к любой точке (нет недостижимых областей).
		
		\item \textbf{(3) Равенство Парсеваля:} Это \textbf{теорема Пифагора} для бесконечности. Квадрат длины вектора равен сумме квадратов его проекций на оси ($c^2 = a^2 + b^2 + \dots$).
		
		\item \textbf{(4) Тривиальный аннулятор:} Если вектор перпендикулярен всем осям координат сразу, то это может быть только нулевой вектор. Если бы нашелся ненулевой перпендикуляр, значит, мы «пропустили» какое-то измерение.
		
	\end{itemize}
	
\end{idea}
	
	\begin{idea}[Логика доказательства]
		
		Самые важные переходы:
		
		\begin{itemize}
			
			\item \textbf{(2) $\Rightarrow$ (1) (От плотности к сходимости ряда):}
			
			Мы знаем, что к вектору $h$ можно приблизиться \textit{какой-то} линейной комбинацией (по определению плотности). Но нам нужно доказать, что именно \textit{ряд Фурье} сходится к $h$.
			
			Здесь работает \textbf{Предложение \ref{minproperty}}: оно гарантирует, что ряд Фурье приближает \textit{лучше всех}. Если какая-то комбинация дает ошибку $\varepsilon$, то ряд Фурье даст ошибку $\le \varepsilon$. Значит, он сходится.
			
			\item \textbf{(1) $\Leftrightarrow$ (3) (Базис $\Leftrightarrow$ Парсеваль):}
			
			Это следует из тождества:
			
			\[ \|\text{Ошибка приближения}\|^2 = \|h\|^2 - \sum |c_k|^2 \]
			
			Если ошибка стремится к 0 (базис), то сумма квадратов стремится к норме (Парсеваль), и наоборот.
			
			
			
			\item \textbf{(2) $\Leftrightarrow$ (4) (Плотность $\Leftrightarrow$ Перпендикуляры):}
			
			Это геометрический факт: пространство $H$ раскладывается на замыкание оболочки $\overline{[e]}$ и её ортогональное дополнение. Чтобы $\overline{[e]}$ было всем пространством, дополнение должно быть нулевым.
			
		\end{itemize}
		
	\end{idea}

\begin{proof}
	\begin{enumerate}
		\item Так как $H$ сепарабельно, в нем существует счетное всюду плотное множество $F = \{f_n\}$ (т.е. $\overline{[F]} = H$).
		\item Применим к $\{f_n\}$ процесс ортогонализации Грама-Шмидта. Получим ортонормированную систему $e = \{e_n\}$.
		\item Процесс сохраняет линейные оболочки: $[e_1, \dots, e_n] = [f_1, \dots, f_n]$. Значит $[e] = [F]$.
		\item Следовательно $\overline{[e]} = \overline{[F]} = H$. По теореме \ref{thm4.4}, система $e$ является базисом.
	\end{enumerate}
\end{proof}

\begin{remark}
	Это означает, что все сепарабельные гильбертовы пространства одной размерности изоморфны (как векторные пространства со скалярным произведением). Например, $L_2[0, 1]$ изоморфно пространству последовательностей $l_2$.
\end{remark}\newpage
	
	% V. ЛИНЕЙНЫЕ ОГРАНИЧЕННЫЕ ОПЕРАТОРЫ
	\section{Линейные ограниченные операторы}
	\hypertarget{bilet12}{}\setupnav{11}{13} % --- Содержимое файла materials/5_1.tex ---

\subsection{Линейные операторы: ограниченность и непрерывность}

Пусть $E_1, E_2$ --- линейные нормированные пространства над полем $\mathbb{K}$ ($\mathbb{R}$ или $\mathbb{C}$).
Отображение $A: E_1 \to E_2$ будем называть \textbf{оператором}, а отображение $f: E \to \mathbb{K}$ --- \textbf{функционалом}.

\begin{definition}[Ограниченный оператор]
	Оператор $A$ называется \textbf{ограниченным}, если он переводит ограниченные множества в ограниченные.
	
	Формально: для любого ограниченного множества $M \subset E_1$ образ $A(M)$ ограничен в $E_2$.
\end{definition}


\begin{idea}
	\textbf{Интуиция:} Ограниченный линейный оператор — это такой оператор, который «не растягивает» векторы бесконечно сильно. У него есть конечный «коэффициент усиления». Если вы подаете на вход вектор длины 1, на выходе не может получиться вектор бесконечной длины.
\end{idea}

\begin{definition}[Образ и Ядро]
	Для линейных операторов вводятся следующие множества:
	\begin{itemize}
		\item \textbf{Образ оператора} (Image):
		\[ \operatorname{Im} A = \{ y \in E_2 \mid \exists x \in E_1 : Ax = y\} \]
		\item \textbf{Ядро оператора} (Kernel):
		\[ \operatorname{Ker} A = \{ x \in E_1 \mid Ax = 0\} \]
	\end{itemize}
\end{definition}


\begin{definition}[Непрерывный оператор]
	Линейный оператор $A$ называется \textbf{непрерывным}, если для любой сходящейся последовательности $x_n \to x$ выполнено $Ax_n \to Ax$.
\end{definition}


\begin{theorem}[Критерий непрерывности линейного оператора]\label{thm:bounded_cont}
	Пусть $E_1, E_2$ --- линейные нормированные пространства, $A: E_1 \to E_2$ --- линейный оператор.
	
	Тогда $A$ ограничен $\iff$ $A$ непрерывен.
\end{theorem}

\begin{proof}
	\textbf{1. Ограниченность $\Rightarrow$ Непрерывность.}
	
	Пусть $A$ ограничен. Воспользуемся эквивалентным определением ограниченности (через норму оператора, см. утверждение \ref{equiv-bounded-oper}): $\exists C > 0: \|Ax\| \leqslant C\|x\|$.
	
	Пусть $x_n \to x$. Это равносильно тому, что $\|x_n - x\|_{E_1} \to 0$. Оценим норму разности образов, используя линейность $A$:
	\[
	\|Ax_n - Ax\|_{E_2} = \|A(x_n - x)\|_{E_2} \leqslant \|A\| \cdot \|x_n - x\|_{E_1}
	\]
	Так как $\|x_n - x\|_{E_1} \to 0$, то и $\|Ax_n - Ax\|_{E_2} \to 0$, что означает $Ax_n \to Ax$.
	
	\textbf{2. Непрерывность $\Rightarrow$ Ограниченность.}
	
	Докажем от противного. Предположим, что $A$ непрерывен, но \textbf{не} является ограниченным.
	
	Это значит, что его «коэффициент растяжения» может быть сколь угодно большим:
	\[ \forall n \in \mathbb{N} \quad \exists x_n \in E_1 : \|Ax_n\|_{E_2} > n \|x_n\|_{E_1} \]
	Очевидно, $x_n \neq 0$. Рассмотрим нормированную и масштабированную последовательность векторов:
	\[ y_n = \frac{x_n}{n \cdot \|x_n\|_{E_1}} \]
	
	Заметим два факта:
	\begin{enumerate}
		\item Последовательность $y_n$ стремится к нулю:
		\[ \|y_n\|_{E_1} = \frac{\|x_n\|}{n \|x_n\|} = \frac{1}{n} \to 0 \implies y_n \to 0 \]
		\item Из непрерывности оператора $A$ следует, что образ должен стремиться к образу нуля:
		\[ A y_n \to A(0) = 0 \]
	\end{enumerate}
	
	Однако, оценим норму образа $y_n$ снизу:
	\[ \|A y_n\|_{E_2} = \left\| A \left( \frac{x_n}{n \|x_n\|} \right) \right\| = \frac{1}{n \|x_n\|} \|A x_n\| > \frac{1}{n \|x_n\|} \cdot (n \|x_n\|) = 1 \]
	
	Получили противоречие: последовательность $A y_n$ должна стремиться к $0$, но норма всех ее элементов больше $1$. Следовательно, предположение неверно, и оператор $A$ ограничен.
\end{proof}

\begin{idea}[Смысл теоремы]
	
	Для линейных операторов понятия «не порвать график» (непрерывность) и «не растягивать бесконечно сильно» (ограниченность) совпадают.
	
	Это очень сильное свойство, которого нет у обычных функций (например, $f(x) = x^2$ непрерывна, но не ограничена на $\mathbb{R}$). 
	
	Причина в \textbf{линейности}: оператор одинаково устроен во всем пространстве. Если он «ведет себя хорошо» (непрерывен) в одной точке (в нуле), он обязан быть ограниченным везде. И наоборот: если у него есть конечный «коэффициент растяжения», он не может совершать скачков.
	
	
	
	2. Идея доказательства 
	
	[Логика доказательства]
	
	\textbf{1. Ограниченность $\Rightarrow$ Непрерывность:}
	
	Здесь работает простая оценка. Если оператор растягивает любой вектор максимум в $\|A\|$ раз, то маленькая разность на входе ($x_n - x \to 0$) превратится в маленькую разность на выходе ($A(x_n - x) \to 0$).
	
	\textbf{2. Непрерывность $\Rightarrow$ Ограниченность (от противного):}
	
	Если оператор \textit{неограничен}, он работает как «усилитель» с бесконечной мощностью: найдутся векторы, которые он удлиняет в $100, 1000, n$ раз.
	
	Мы используем линейность, чтобы «обмануть» его: берем эти векторы и уменьшаем их длину до микроскопических размеров ($y_n \to 0$). Но так как коэффициент растяжения растет быстрее ($n$), на выходе мы все равно получаем «большие» векторы ($\|Ay_n\| > 1$). Это противоречит непрерывности в нуле.
	
\end{idea}\newpage
	\hypertarget{bilet13}{}\setupnav{12}{14} 
\subsection{Топологии и сходимости в пространстве операторов L(E1,E2).Норма оператора. Полнота нормированного пространства  L(E1,E2) (в случае, когда E2 — банахово)}

\begin{theorem}[Эквивалентные определения ограниченности]\label{equiv-bounded-oper}
	Пусть $A: E_1 \to E_2$ --- линейный оператор. Следующие условия эквивалентны:
	\begin{enumerate}
		\item $A$ --- ограниченный (переводит ограниченные множества в ограниченные);
		\item $\exists K \geqslant 0: \forall x \in E_1 \quad \|Ax\|_{E_2} \leqslant K \|x\|_{E_1}$;
		\item Образ единичного шара $B = \{x : \|x\| \leqslant 1\}$ ограничен.
	\end{enumerate}
\end{theorem}

\begin{definition}[Норма оператора]
	\textbf{Нормой} линейного ограниченного оператора $A$ называется точная нижняя грань констант, ограничивающих оператор:
	\[ \|A\| = \inf \{ K : \forall x \in E_1, \|Ax\| \leqslant K \|x\| \} \]
\end{definition}

\begin{idea}
	Норма оператора $\|A\|$ — это «максимальный коэффициент растяжения». Это число показывает, во сколько раз оператор может увеличить длину вектора в худшем случае.
\end{idea}


\begin{proposition}[Вычисление нормы]
	Пусть $A$ --- линейный ограниченный оператор. Тогда:
	\[ \|A\| = \sup_{\|x\| \leqslant 1} \|Ax\| = \sup_{\|x\| = 1} \|Ax\| = \sup_{x \neq 0} \frac{\|Ax\|}{\|x\|} \]
\end{proposition}

\begin{definition}
	$\mathcal{L}(E_1, E_2)$ --- множество всех линейных ограниченных операторов из $E_1$ в $E_2$. Оно является линейным пространством с введенной выше нормой $\|A\|$.
\end{definition}

\begin{definition}[Сопряженное пространство]
	Пространство линейных непрерывных функционалов на $E$ называется \textbf{сопряженным (или двойственным)} пространством:
	\[ E^* = \mathcal{L}(E, \mathbb{K}) \]
\end{definition}

\begin{theorem}[О полноте пространства операторов]
	Пусть $E_1, E_2$ --- линейные нормированные пространства.
	\begin{enumerate}
		\item $\mathcal{L}(E_1, E_2)$ всегда является линейным нормированным пространством.
		\item Если $E_2$ --- \textbf{банахово} (полное), то $\mathcal{L}(E_1, E_2)$ тоже \textbf{банахово}.
	\end{enumerate}
\end{theorem}

\begin{remark}
	Полнота $E_1$ не требуется! Важна только полнота пространства, куда мы отображаем.
\end{remark}

\begin{proof}
	\textbf{1. Аксиомы нормы.}
	Линейность очевидна. Проверим неравенство треугольника:
	\[ \|A+B\| = \sup_{\|x\|=1} \|Ax + Bx\| \leqslant \sup_{\|x\|=1} (\|Ax\| + \|Bx\|) \leqslant \sup \|Ax\| + \sup \|Bx\| = \|A\| + \|B\|. \]
	
	\textbf{2. Полнота.}
	Пусть $E_2$ полно. Рассмотрим фундаментальную последовательность операторов $\{A_n\} \subset \mathcal{L}$.
	\[ \forall \varepsilon > 0 \exists N: \forall n, m \geqslant N \implies \|A_n - A_m\| < \varepsilon. \]
	
	\textbf{Шаг 1: Поточечная сходимость.}
	Зафиксируем произвольный $x \in E_1$.
	\[ \|A_n x - A_m x\| \leqslant \|A_n - A_m\| \cdot \|x\| \to 0. \]
	Значит, последовательность векторов $\{A_n x\}$ фундаментальна в $E_2$.
	Так как $E_2$ банахово, существует предел. Обозначим его $Ax := \lim_{n \to \infty} A_n x$.
	Оператор $A$ линеен (предел линейных отображений).
	
	\textbf{Шаг 2: Ограниченность предельного оператора $A$.}
	Фундаментальная последовательность $\{A_n\}$ ограничена по норме: $\exists C: \|A_n\| \leqslant C$.
	Тогда для любого $x$: $\|A_n x\| \leqslant C \|x\|$. Переходя к пределу: $\|Ax\| \leqslant C \|x\|$.
	Значит, $A \in \mathcal{L}(E_1, E_2)$.
	
	\textbf{Шаг 3: Сходимость по норме оператора ($\|A_n - A\| \to 0$).}
	В условии фундаментальности $\|A_n x - A_m x\| \leqslant \varepsilon \|x\|$ перейдем к пределу при $m \to \infty$:
	\[ \|A_n x - Ax\| \leqslant \varepsilon \|x\| \quad (\forall n \geqslant N, \forall x). \]
	Это означает, что $\sup_{\|x\|=1} \|(A_n - A)x\| \leqslant \varepsilon$, то есть $\|A_n - A\| \leqslant \varepsilon$.
	Следовательно, $A_n \to A$ в пространстве $\mathcal{L}$.
\end{proof}

\begin{corollary}
	Пространство $E^*$ всегда полное (так как $\mathbb{K} = \mathbb{R}$ или $\mathbb{C}$ полны).
\end{corollary}\newpage
	\hypertarget{bilet14}{}\setupnav{13}{15} % --- Содержимое файла materials/5_3.tex ---

\subsection{Продолжение оператора по непрерывности}

\begin{idea}[Смысл теоремы]
	Часто бывает сложно задать оператор сразу на всем «сложном» пространстве (например, $L_2$). Гораздо проще задать его на «хорошем» плотном подмножестве (например, на многочленах или непрерывных функциях), а затем «растянуть» на все остальное пространство с помощью предельного перехода.
	
	Эта теорема гарантирует, что такое «растягивание» возможно и корректно, \textbf{только если} пространство, куда мы попадаем ($E_2$), является \textbf{полным} (Банаховым). Иначе пределов последовательностей может просто не существовать.
\end{idea}

\begin{theorem}[О продолжении по непрерывности]\label{extrapolate-to-closure}
	Пусть:
	\begin{itemize}
		\item $E_1$ --- линейное нормированное пространство;
		\item $E_2$ --- \textbf{банахово} пространство (полное);
		\item $D(A) \subset E_1$ --- линейное подпространство, всюду плотное в $E_1$ ($\overline{D(A)} = E_1$);
		\item $A: D(A) \to E_2$ --- линейный ограниченный оператор.
	\end{itemize}
	
	Тогда существует \textbf{единственное} продолжение оператора $A$ до линейного непрерывного оператора $\widetilde{A}$, определенного на всем $E_1$, причем их нормы совпадают:
	\[ \exists! \widetilde{A} \in \mathcal{L}(E_1, E_2): \quad \widetilde{A}\big|_{D(A)} = A \quad \text{и} \quad \|\widetilde{A}\| = \|A\|. \]
\end{theorem}

\begin{proof}
	\textbf{1. Единственность.}
	Пусть существуют два таких непрерывных расширения $\widetilde{A}_1$ и $\widetilde{A}_2$.
	Возьмем произвольный $x \in E_1$. В силу плотности $D(A)$, существует последовательность $\{x_n\} \subset D(A)$, сходящаяся к $x$.
	В силу непрерывности операторов:
	\[ \widetilde{A}_1 x = \lim_{n \to \infty} \widetilde{A}_1 x_n \overset{x_n \in D(A)}{=} \lim_{n \to \infty} A x_n = \lim_{n \to \infty} \widetilde{A}_2 x_n = \widetilde{A}_2 x. \]
	Значит, $\widetilde{A}_1 = \widetilde{A}_2$.
	
	\textbf{2. Существование (Конструкция).}
	Для любого $x \in E_1$ выберем последовательность $\{x_n\} \subset D(A)$, такую что $x_n \to x$.
	Определим значение оператора формулой:
	\[ \widetilde{A}x := \lim_{n \to \infty} A x_n \]
	
	Чтобы это определение было корректным, нужно проверить три факта:
	
	\textbf{а) Предел существует.}
	Так как $x_n \to x$, последовательность $\{x_n\}$ фундаментальна в $E_1$.
	В силу ограниченности (непрерывности) исходного оператора $A$:
	\[ \|Ax_n - Ax_m\| = \|A(x_n - x_m)\| \leqslant \|A\| \cdot \|x_n - x_m\| \to 0. \]
	Значит, последовательность образов $\{Ax_n\}$ фундаментальна в $E_2$.
	Так как $E_2$ --- \textbf{банахово}, эта последовательность имеет предел.
	
	
	
	\textbf{б) Предел не зависит от выбора последовательности.}
	Пусть $x_n \to x$ и $x'_n \to x$. Рассмотрим смешанную последовательность $y = \{x_1, x'_1, x_2, x'_2, \dots\}$.
	Очевидно, $y_k \to x$. Согласно пункту (а), последовательность $\{Ay_k\}$ сходится, а значит, все её подпоследовательности (в частности $\{Ax_n\}$ и $\{Ax'_n\}$) сходятся к одному и тому же пределу.
	
	\textbf{в) Совпадение на $D(A)$.}
	Если $x \in D(A)$, то в качестве последовательности можно взять стационарную $x_n = x$.
	Тогда $\widetilde{A}x = \lim Ax = Ax$.
	
	\textbf{3. Линейность.}
	Следует из линейности $A$ и свойств пределов:
	\[ \widetilde{A}(\alpha x + \beta y) = \lim A(\alpha x_n + \beta y_n) = \alpha \lim Ax_n + \beta \lim Ay_n = \alpha \widetilde{A}x + \beta \widetilde{A}y. \]
	
	\textbf{4. Сохранение нормы.}
	\begin{itemize}
		\item Оценка сверху ($\|\widetilde{A}\| \leqslant \|A\|$):
		\[ \|\widetilde{A}x\| = \lim_{n \to \infty} \|Ax_n\| \leqslant \lim_{n \to \infty} \|A\| \cdot \|x_n\| = \|A\| \cdot \|x\|. \]
		Следовательно, $\|\widetilde{A}\| \leqslant \|A\|$.
		\item Оценка снизу ($\|\widetilde{A}\| \geqslant \|A\|$):
		Так как $\widetilde{A}$ является расширением $A$, то супремум на всем пространстве не может быть меньше супремума на подмножестве:
		\[ \|\widetilde{A}\| = \sup_{x \in E_1, \|x\|=1} \|\widetilde{A}x\| \geqslant \sup_{x \in D(A), \|x\|=1} \|\widetilde{A}x\| = \sup_{x \in D(A), \|x\|=1} \|Ax\| = \|A\|. \]
	\end{itemize}
	Итого: $\|\widetilde{A}\| = \|A\|$.
\end{proof}\newpage
	\hypertarget{bilet15}{}\setupnav{14}{16} % --- Содержимое файла materials/5_4.tex ---

\subsection{Теорема Банаха--Штейнгауза}

\begin{idea}[Смысл теоремы]
	\textbf{Принцип равномерной ограниченности.}
	Представьте, что у вас есть набор пружин (операторов). Вы проверяете их по одной: подвешиваете груз весом 1 кг, и ни одна пружина не растягивается до бесконечности. Подвешиваете 2 кг — тоже все держат.
	
	Теорема утверждает: если пружины «держат» любой конкретный вес (поточечная ограниченность), то существует \textbf{общий предел жесткости} для всех пружин сразу (равномерная ограниченность). Не может быть такого, что пружины становятся всё слабее и слабее, даже если ни одна не ломается полностью.
	
	\textbf{Важно:} Это работает только если пространство $X$ «жирное» (полное). В «дырявых» пространствах теорема ломается.
\end{idea}

\begin{theorem}[Банаха--Штейнгауза]\label{Banach-Steingaus}
	Пусть $X$ --- банахово пространство (полное линейное нормированное), $Y$ --- линейное нормированное пространство.
	Пусть $\mathcal{A} \subseteq \mathcal{L}(X, Y)$ --- семейство линейных непрерывных операторов.
	
	Если это семейство ограничено в каждой точке:
	\[ \forall x \in X \quad \sup_{A \in \mathcal{A}} \|Ax\|_Y < \infty, \]
	то оно ограничено равномерно (по норме операторов):
	\[ \sup_{A \in \mathcal{A}} \|A\| < \infty. \]
\end{theorem}

\begin{idea}[Идея доказательства]
	Мы разбиваем пространство $X$ на множества $X_n$ («области», где операторы не растягивают векторы сильнее, чем в $n$ раз).
	Так как для любой точки $x$ операторы ограничены, объединение всех $X_n$ дает всё пространство $X$.
	
	Здесь вступает в игру **Теорема Бэра о категориях**: полное пространство нельзя составить из «тощих» (нигде не плотных) множеств. Значит, хотя бы одно множество $X_m$ должно быть «жирным» — содержать внутри себя целый шар.
	Используя линейность, мы «разносим» ограниченность с этого шара на всё пространство.
\end{idea}


\begin{proof}
	\textbf{1. Построение множеств Лебега.}
	Рассмотрим множества, на которых семейство операторов ограничено числом $n$:
	\[ X_n = \left\{x \in X \mid \sup_{A \in \mathcal{A}} \|Ax\|_Y \leqslant n \right\}. \]
	Условие $\sup \|Ax\| \leqslant n$ эквивалентно тому, что $\forall A \in \mathcal{A}: \|Ax\| \leqslant n$. Поэтому:
	\[ X_n = \bigcap_{A \in \mathcal{A}} \{x \in X \mid \|Ax\|_Y \leqslant n\}. \]
	
	\textbf{2. Замкнутость $X_n$.}
	Для фиксированного $A$ множество $\{x \mid \|Ax\| \leqslant n\}$ является прообразом замкнутого шара $\overline{B}_n(0) \subset Y$ при непрерывном отображении $A$. Следовательно, оно замкнуто.
	Множество $X_n$ есть пересечение замкнутых множеств, значит, оно тоже \textbf{замкнуто}.
	
	\textbf{3. Применение теоремы Бэра.}
	По условию теоремы (поточечная ограниченность), для каждого $x$ супремум конечен, значит, каждый $x$ попадает в какое-то $X_n$.
	\[ X = \bigcup_{n=1}^{\infty} X_n. \]
	Так как $X$ — полное пространство, по \textbf{теореме Бэра} оно не может быть объединением счетного числа нигде не плотных множеств. Значит, существует номер $m$, для которого $X_m$ имеет внутреннюю точку.
	То есть существуют $x_0 \in X$ и $\varepsilon > 0$, такие что шар лежит в множестве:
	\[ \overline{B}_\varepsilon(x_0) \subset X_m. \]
	Это означает, что для любого $z \in \overline{B}_\varepsilon(x_0)$ и любого $A \in \mathcal{A}$ выполняется $\|Az\| \leqslant m$.
	
	\textbf{4. Оценка нормы операторов.}
	Возьмем произвольный $u \in X$ с единичной нормой ($\|u\|=1$).
	Представим его через векторы из шара: $x_0$ и $x_0 + \varepsilon u$.
	\[ A(u) = A\left( \frac{(x_0 + \varepsilon u) - x_0}{\varepsilon} \right) = \frac{1}{\varepsilon} (A(x_0 + \varepsilon u) - A(x_0)). \]
	Оценим норму, используя неравенство треугольника:
	\[ \|Au\| \leqslant \frac{1}{\varepsilon} \left( \|A(x_0 + \varepsilon u)\| + \|A(x_0)\| \right). \]
	Так как $x_0 \in X_m$ и $(x_0 + \varepsilon u) \in \overline{B}_\varepsilon(x_0) \subset X_m$, то нормы их образов не превосходят $m$:
	\[ \|Au\| \leqslant \frac{1}{\varepsilon} (m + m) = \frac{2m}{\varepsilon}. \]
	
	Эта оценка не зависит ни от $u$, ни от выбора оператора $A$.
	Следовательно, $\sup_{A \in \mathcal{A}} \|A\| \leqslant \frac{2m}{\varepsilon} < \infty$.
\end{proof}

\begin{corollary}[О предельном операторе]
	Пусть $X$ — банахово пространство, $Y$ — линейное нормированное.
	Пусть $\{A_n\} \subset \mathcal{L}(X, Y)$ — последовательность непрерывных операторов, которая сходится поточечно:
	\[ \forall x \in X \quad A_n x \to A x. \]
	Тогда предельный оператор $A$ также является \textbf{линейным и непрерывным} ($A \in \mathcal{L}(X, Y)$).
\end{corollary}

\begin{proof}
	Линейность предела очевидна. Докажем непрерывность (ограниченность).
	Так как последовательность $\{A_n x\}$ сходится, она ограничена для каждого $x$.
	По теореме Банаха-Штейнгауза, нормы $\|A_n\|$ ограничены в совокупности константой $C$.
	Тогда для любого $x$:
	\[ \|Ax\| = \lim_{n \to \infty} \|A_n x\| \leqslant \limsup_{n \to \infty} (\|A_n\| \cdot \|x\|) \leqslant C \|x\|. \]
	Значит, $A$ ограничен.
\end{proof}

\begin{remark}
	Теоремы о сходимости квадратурных формул и интерполяционных процессов (Билет №16) являются прямыми следствиями теоремы Банаха--Штейнгауза.
\end{remark}\newpage
	\hypertarget{bilet16}{}\setupnav{15}{17} % --- Содержимое файла materials/5_5.tex ---

\subsection{Поточечная сходимость последовательностей операторов}

\begin{idea}[Смысл теоремы]
	Если последовательность операторов сходится в каждой точке, то предельный оператор «наследует» линейность и непрерывность.
	\\
	\textbf{Важно:} Это работает, только если пространство $E_1$ полное (чтобы работала теорема Банаха-Штейнгауза для оценки норм) и $E_2$ полное (чтобы предел вообще существовал).
\end{idea}

\begin{theorem}[О полноте относительно поточечной сходимости]
	Пусть $E_1, E_2$ --- банаховы пространства.
	Пусть $\{A_n\} \subset \mathcal{L}(E_1, E_2)$ --- последовательность операторов, которая является фундаментальной в каждой точке:
	\[ \forall x \in E_1 \quad \{A_n x\} \text{ --- фундаментальна в } E_2. \]
	Тогда существует предельный линейный непрерывный оператор $A \in \mathcal{L}(E_1, E_2)$ такой, что:
	\[ \forall x \in E_1 \quad \lim_{n \to \infty} A_n x = Ax. \]
\end{theorem}

\begin{proof}
	\textbf{1. Существование и Линейность.}
	Так как для любого $x$ последовательность $\{A_n x\}$ фундаментальна, а пространство $E_2$ полное (банахово), то предел существует. Определим оператор $A$:
	\[ Ax := \lim_{n \to \infty} A_n x. \]
	Линейность $A$ следует из свойств предела и линейности каждого $A_n$.
	
	\textbf{2. Непрерывность (Ограниченность).}
	Так как последовательность $\{A_n x\}$ сходится, она ограничена для каждого $x$.
	Значит, семейство $\{A_n\}$ поточечно ограничено.
	Так как $E_1$ --- банахово, по теореме \textbf{Банаха--Штейнгауза} (Thm \ref{Banach-Steingaus}) оно ограничено равномерно:
	\[ \exists M > 0: \quad \forall n \in \mathbb{N} \quad \|A_n\| \leqslant M. \]
	
	Оценим норму предельного оператора. Для любого $x$:
	\[ \|A_n x\| \leqslant \|A_n\| \cdot \|x\| \leqslant M \|x\|. \]
	Переходя к пределу при $n \to \infty$ (учитывая непрерывность нормы):
	\[ \|Ax\| \leqslant M \|x\|. \]
	Следовательно, $A$ ограничен и $A \in \mathcal{L}(E_1, E_2)$.
\end{proof}

\begin{idea}[Метод плотности]
	Как проверить, что операторы $A_n$ сходятся к $A$ во всём бесконечном пространстве? Проверять каждую точку невозможно.
	Критерий ниже дает алгоритм:
	1. Проверь \textbf{стабильность}: нормы операторов не должны расти бесконечно ($\|A_n\| \le M$).
	2. Проверь \textbf{сходимость на простых векторах} (на базисе, многочленах и т.д.).
	Если оба условия выполнены — сходимость есть везде.
\end{idea}


\begin{theorem}[Критерий поточечной сходимости]\label{bshetengcorollary}
	Пусть $E_1$ --- банахово, $E_2$ --- линейное нормированное пространство.
	Пусть $\{A_n\} \subset \mathcal{L}(E_1, E_2)$ и $A \in \mathcal{L}(E_1, E_2)$.
	
	Последовательность $A_n$ сходится к $A$ поточечно ($\forall x: A_n x \to Ax$) тогда и только тогда, когда выполнены два условия:
	\begin{enumerate}
		\item Последовательность норм ограничена: $\exists M: \sup_n \|A_n\| \leqslant M$.
		\item Сходимость есть на плотном множестве: $\exists S \subset E_1$, $\overline{[S]} = E_1$ такое, что $\forall s \in S: A_n s \to As$.
	\end{enumerate}
\end{theorem}

\begin{proof}
	\textbf{Необходимость ($\Rightarrow$)}
	\begin{itemize}
		\item Сходимость на $S$ следует из сходимости на всем $E_1$ (так как $S \subset E_1$).
		\item Ограниченность норм следует из теоремы Банаха--Штейнгауза (так как есть поточечная сходимость на всем $E_1$, значит есть и поточечная ограниченность).
	\end{itemize}
	
	\textbf{Достаточность ($\Leftarrow$)}
	Пусть условия 1 и 2 выполнены. Зафиксируем произвольный $x \in E_1$ и $\varepsilon > 0$.
	Так как линейная оболочка $[S]$ плотна в $E_1$, найдем элемент $y \in [S]$, близкий к $x$:
	\[ \|x - y\| < \frac{\varepsilon}{3M^*}, \quad \text{где } M^* = \max(M, \|A\|). \]
	
	Оценим разность $\|A_n x - Ax\|$ методом «$\varepsilon/3$»:
	\[ \|A_n x - Ax\| \leqslant \underbrace{\|A_n x - A_n y\|}_{\text{(1)}} + \underbrace{\|A_n y - Ay\|}_{\text{(2)}} + \underbrace{\|Ay - Ax\|}_{\text{(3)}}. \]
	
	Разберем слагаемые:
	\begin{itemize}
		\item (1) $\|A_n(x - y)\| \leqslant \|A_n\| \|x - y\| \leqslant M \cdot \frac{\varepsilon}{3M^*} \leqslant \frac{\varepsilon}{3}$.
		\item (3) $\|A(y - x)\| \leqslant \|A\| \|x - y\| \leqslant \|A\| \cdot \frac{\varepsilon}{3M^*} \leqslant \frac{\varepsilon}{3}$.
		\item (2) Так как $y \in [S]$, то $A_n y \to Ay$. Значит, $\exists N$, что при $n > N$: $\|A_n y - Ay\| < \frac{\varepsilon}{3}$.
	\end{itemize}
	
	Суммируя, получаем:
	\[ \|A_n x - Ax\| < \frac{\varepsilon}{3} + \frac{\varepsilon}{3} + \frac{\varepsilon}{3} = \varepsilon. \]
	Значит, $A_n x \to Ax$ для любого $x$.
\end{proof}

\begin{idea}
	Этот критерий лежит в основе обоснования многих численных методов (например, сходимости квадратурных формул). Мы проверяем сходимость для полиномов (это легко), убеждаемся в ограниченности норм сумм, и получаем сходимость для любой непрерывной функции.
\end{idea}\newpage
	
	% VI. СОПРЯЖЕННОЕ ПРОСТРАНСТВО
	\section{Сопряженное пространство}
	\hypertarget{bilet17}{}\setupnav{16}{18} % --- Содержимое файла materials/6_1.tex ---

\subsection{Теорема Рисса--Фреше}

\begin{definition}[Линейный функционал]
	Линейный оператор, действующий из нормированного пространства $X$ в скалярное поле $\mathbb{K}$ (где $\mathbb{K} = \mathbb{R}$ или $\mathbb{C}$), называется \textbf{линейным функционалом}.
	
	Множество всех линейных непрерывных функционалов над $X$ образует \textbf{сопряженное пространство}:
	\[ X^* = \mathcal{L}(X, \mathbb{K}) \]
\end{definition}

\begin{idea}[Смысл теоремы Рисса--Фреше]
	\textbf{Функционал = Вектор.}
	В «хороших» (гильбертовых) пространствах любой инструмент измерения (функционал) сам является вектором этого пространства.
	
	Измерение длины проекции вектора $x$ на вектор $y$ записывается как скалярное произведение $(y, x)$. Теорема утверждает, что \textbf{любой} линейный непрерывный функционал $f(x)$ выглядит именно так. Нет никаких других «экзотических» функционалов.
	
	Геометрически вектор $y$ — это градиент функционала, перпендикуляр к его ядру (нулевому уровню).
\end{idea}



\begin{theorem}[Рисса--Фреше]\label{Riss-Freshe}
	Пусть $H$ — гильбертово пространство, $f \in H^*$ — линейный ограниченный функционал.
	Тогда существует \textbf{единственный} вектор $y \in H$ такой, что функционал представим в виде скалярного произведения:
	\[ \forall x \in H: \quad f(x) = \langle y, x \rangle. \]
	При этом их нормы совпадают: $\|f\|_{H^*} = \|y\|_H$.
\end{theorem}

\begin{proof}
	\textbf{1. Существование.}
	
	\begin{itemize}
		
		\item \textit{Случай $f \equiv 0$.}
		
		Берём $y=0$. Тогда для любого $x\in H$:
		$$
		f(x)=0=\langle 0,x\rangle,
		$$
		и $\|f\|=\|y\|=0$.
		
		Далее считаем $f\not\equiv 0$.
		
		\medskip
		
		\item \textbf{Шаг 1. Рассматриваем ядро функционала.}
		
		Обозначим
		$$
		M:=\ker f=\{x\in H:\ f(x)=0\}.
		$$
		
		Тогда:
		\begin{itemize}
			\item $M\neq H$ (так как $f\neq 0$);
			\item $M$ замкнуто, поскольку $f$ непрерывен.
		\end{itemize}
		
		По теореме о разложении гильбертова пространства:
		$$
		H = M \oplus M^\perp,
		$$
		поэтому существует вектор
		$$
		b\in M^\perp,\quad b\neq 0.
		$$
		
		\medskip
		
		\item \textbf{Шаг 2. Стратегия.}
		
		Мы хотим получить формулу вида
		$$
		f(x)=C\langle b,x\rangle.
		$$
		
		Так как $f=0$ на $M$, естественно попытаться для каждого $x$
		вычесть из него компоненту вдоль $b$, чтобы попасть в ядро.
		
		То есть ищем число $\alpha$, для которого
		$$
		x-\alpha b\in M=\ker f.
		$$
		
		\medskip
		
		\item \textbf{Шаг 3. Подбор коэффициента.}
		
		Требуем:
		$$
		f(x-\alpha b)=0.
		$$
		
		По линейности:
		$$
		f(x-\alpha b)=f(x)-\alpha f(b).
		$$
		
		Так как $b\notin \ker f$ (иначе $b\perp b\Rightarrow b=0$),
		имеем $f(b)\neq 0$. Поэтому
		$$
		f(x)-\alpha f(b)=0
		\quad\Rightarrow\quad
		\alpha=\frac{f(x)}{f(b)}.
		$$
		
		Положим
		$$
		p_x:=x-\frac{f(x)}{f(b)}\,b.
		$$
		Тогда $p_x\in\ker f$.
		
		\medskip
		
		\item \textbf{Шаг 4. Используем ортогональность.}
		
		Так как $b\perp\ker f$, получаем:
		$$
		\langle b,p_x\rangle=0.
		$$
		
		Подставляя $p_x$:
		$$
		0=\left\langle b,\;x-\frac{f(x)}{f(b)}b\right\rangle
		=\langle b,x\rangle-\frac{f(x)}{f(b)}\|b\|^2.
		$$
		
		\medskip
		
		\item \textbf{Шаг 5. Выражаем функционал.}
		
		Из последнего равенства:
		$$
		f(x)=\frac{f(b)}{\|b\|^2}\,\langle b,x\rangle.
		$$
		
		Перенеся скаляр внутрь скалярного произведения (в первое место), получаем
		$$
		f(x)=\left\langle
		\frac{\overline{f(b)}}{\|b\|^2}\,b,\;x
		\right\rangle.
		$$
		
		Обозначим
		$$
		y:=\frac{\overline{f(b)}}{\|b\|^2}\,b.
		$$
		
		Тогда для всех $x\in H$:
		$$
		f(x)=\langle y,x\rangle.
		$$
		
		Существование доказано.
		
	\end{itemize}
	
	\medskip
	
	\textbf{2. Единственность.}
	
	Пусть
	$$
	\langle y,x\rangle=\langle z,x\rangle
	\quad \forall x\in H.
	$$
	
	Тогда
	$$
	\langle y-z,x\rangle=0\quad \forall x.
	$$
	
	Подставляя $x=y-z$, получаем:
	$$
	\|y-z\|^2=0 \;\Rightarrow\; y=z.
	$$
	
	\medskip
	
	\textbf{3. Равенство норм.}
	
	\begin{itemize}
		
		\item \textit{Оценка сверху.}
		
		По неравенству Коши–Буняковского:
		$$
		|f(x)|=|\langle y,x\rangle|
		\le \|y\|\,\|x\|.
		$$
		
		Следовательно,
		$$
		\|f\|=\sup_{\|x\|=1}|f(x)|
		\le \|y\|.
		$$
		
		\medskip
		
		\item \textit{Оценка снизу.}
		
		Если $y\neq 0$, берём $x=y$:
		$$
		f(y)=\langle y,y\rangle=\|y\|^2.
		$$
		
		По определению нормы функционала:
		$$
		|f(y)|\le \|f\|\,\|y\|.
		$$
		
		Отсюда:
		$$
		\|y\|^2\le \|f\|\,\|y\|
		\;\Rightarrow\;
		\|y\|\le \|f\|.
		$$
		
		(При $y=0$ равенство очевидно.)
		
	\end{itemize}
	
	Из двух неравенств:
	$$
	\|f\|=\|y\|.
	$$
	
\end{proof}
\newpage
	\hypertarget{bilet18}{}\setupnav{17}{19} % --- Содержимое файла materials/6_2.tex ---

\subsection{Теорема Хана--Банаха}

\begin{idea}[Смысл теоремы]
	\textbf{Принцип сохранения нормы.}
	Представьте, что вы измеряете "длину" векторов только на плоскости внутри трехмерного пространства. У вас есть линейка (функционал).
	Теорема Хана-Банаха говорит, что вы можете использовать эту же линейку для измерения любых векторов в пространстве, при этом:
	1. Результаты измерений на плоскости не изменятся.
	2. "Масштаб" линейки (её норма) не увеличится.
	
	Это позволяет нам строить функционалы с нужными свойствами "по кирпичикам", начиная с малых подпространств.
\end{idea}

\begin{theorem}[Хана--Банаха]\label{Han-Banah}
	Пусть $E$ --- линейное нормированное пространство, $M \subset E$ --- линейное многообразие.
	Пусть $f$ --- линейный ограниченный функционал, заданный на $M$.
	
	Тогда существует функционал $\widetilde{f} \in E^*$ (определенный на всем $E$), который является продолжением $f$ с сохранением нормы:
	\begin{enumerate}
		\item $\widetilde{f}\big|_{M} = f$;
		\item $\|\widetilde{f}\|_{E^*} = \|f\|_{M^*}$.
	\end{enumerate}
\end{theorem}

\begin{proof}[Доказательство для вещественного сепарабельного случая]
	
	\textbf{Шаг 1. "Завоевание одного измерения".}
	Пусть $M \neq E$. Возьмем вектор $x_0 \notin M$. Рассмотрим подпространство
	$$
	M_1 = M \oplus [x_0].
	$$
	Любой элемент $y \in M_1$ однозначно представим в виде
	$$
	y = x + \alpha x_0,\quad x \in M.
	$$
	
	Зададим продолжение $f_1$ формулой:
	$$
	f_1(y) = f(x) + \alpha\, c,
	$$
	где
	$$
	c := f_1(x_0)
	$$
	— пока неизвестное число.
	
	Мы хотим выбрать $c$ так, чтобы норма не выросла:
	$$
	\|f_1\| \leqslant \|f\|
	$$
	(обратное неравенство очевидно). То есть для любого $y$:
	$$
	|f(x) + \alpha c| \leqslant \|f\| \cdot \|x + \alpha x_0\|.
	$$
	
	Если $\alpha = 0$, это верно. Если $\alpha \neq 0$, поделим на $|\alpha|$.
	Пусть $z = x/\alpha$. Неравенство примет вид:
	$$
	|f(z) + c| \leqslant \|f\| \cdot \|z + x_0\|.
	$$
	
	Раскроем модуль (двойное неравенство):
	$$
	-\|f\|\|z+x_0\| - f(z) \leqslant c \leqslant \|f\|\|z+x_0\| - f(z).
	$$
	
	Чтобы такое число $c$ существовало, нужно, чтобы \textbf{левая часть}
	(для любого $z_1$) была меньше или равна \textbf{правой части}
	(для любого $z_2$):
	$$
	-f(z_1) - \|f\|\|z_1+x_0\|
	\leqslant
	-f(z_2) + \|f\|\|z_2+x_0\|.
	$$
	
	Перегруппируем:
	$$
	f(z_2) - f(z_1)
	\leqslant
	\|f\|\bigl(\|z_1+x_0\| + \|z_2+x_0\|\bigr).
	$$
	
	Это действительно верно, так как:
	$$
	f(z_2) - f(z_1) = f(z_2 - z_1)
	\leqslant \|f\|\|z_2 - z_1\|
	= \|f\|\|(z_2+x_0) - (z_1+x_0)\|
	\leqslant \|f\|\bigl(\|z_2+x_0\| + \|z_1+x_0\|\bigr).
	$$
	
	Следовательно, существует число $c$, удовлетворяющее условию.
	Мы расширили функционал на одно измерение с сохранением нормы.
	
	\medskip
	\textbf{Шаг 2. Индукция (для сепарабельного случая).}
	
	Пусть $E$ сепарабельно. Выберем счетное всюду плотное множество $\{x_n\}$.
	
	Построим последовательность вложенных подпространств
	$$
	M_0 = M,\qquad
	M_n = M_{n-1} \oplus [x_{n-1}]
	$$
	(если $x_{n-1} \notin M_{n-1}$).
	
	На каждом шаге применяем Шаг 1 и получаем функционал $f_n$ с той же нормой.
	
	Определим функционал $f_\infty$ на объединении
	$$
	M_\infty = \bigcup M_n.
	$$
	Это линейное многообразие плотно в $E$.
	
	Так как функционал ограничен, по теореме о продолжении по непрерывности
	(Thm \ref{extrapolate-to-closure}) его можно единственным образом
	продолжить на всё $E$ с сохранением нормы.
\end{proof}

\subsubsection*{Следствия из Теоремы Хана-Банаха}

\begin{corollary}[О разделении точки и подпространства]\label{cor:HB_sep}
	Пусть $M \subsetneq E$ --- линейное подпространство, $x_0 \notin \overline{M}$.
	Тогда существует функционал $f \in E^*$, такой что:
	\begin{enumerate}
		\item $f(x) = 0$ для всех $x \in M$ ($M \subset \ker f$);
		\item $f(x_0) = 1$;
		\item $\|f\| = \frac{1}{\rho(x_0, M)}$.
	\end{enumerate}
\end{corollary}

\begin{idea}
	Это следствие говорит, что если точка не лежит в подпространстве (или его замыкании), то всегда найдется функционал, который "видит" эту разницу: он обнуляет всё подпространство, но на точке выдает единицу.
\end{idea}


\begin{proof}
	Рассмотрим $M_1 = M \oplus [x_0]$. Любой $y \in M_1$ имеет вид $y = x + \alpha x_0$.
	Зададим функционал $f_1$ на $M_1$: $f_1(x + \alpha x_0) = \alpha$.
	Очевидно, $f_1|_M = 0$ и $f_1(x_0) = 1$.
	
	Найдем норму $f_1$.
	\[ \|f_1\| = \sup_{y \neq 0} \frac{|f_1(y)|}{\|y\|} = \sup_{\alpha \neq 0, x \in M} \frac{|\alpha|}{\|x + \alpha x_0\|} = \sup \frac{1}{\|x/\alpha + x_0\|} = \frac{1}{\inf_{z \in M} \|z + x_0\|} = \frac{1}{\rho(x_0, M)}. \]
	По теореме Хана-Банаха продолжаем $f_1$ на всё $E$ с сохранением этой нормы.
\end{proof}

\begin{corollary}[Существование опорного функционала]\label{cor:HB_norm}
	Для любого $x \in E$ существует функционал $f \in E^*$, такой что:
	\begin{enumerate}
		\item $\|f\| = 1$;
		\item $f(x) = \|x\|$.
	\end{enumerate}
\end{corollary}

\begin{proof}
	Применим предыдущее следствие для $M = \{0\}$ и точки $x_0 = x$.
	Тогда $\rho(x, \{0\}) = \|x\|$.
	Существует $g \in E^*$ с нормой $1/\|x\|$, такой что $g(x) = 1$.
	Положим $f = \|x\| \cdot g$. Тогда $\|f\| = \|x\| \cdot (1/\|x\|) = 1$ и $f(x) = \|x\| \cdot 1 = \|x\|$.
\end{proof}

\begin{corollary}[Критерий равенства нулю]\label{hbcorollary3}
	Следующие утверждения эквивалентны:
	\begin{enumerate}
		\item $x = 0$;
		\item $\forall f \in E^*: f(x) = 0$.
	\end{enumerate}
\end{corollary}
\begin{proof}
	Если $x \neq 0$, то по следствию \ref{cor:HB_norm} существует $f$ с $f(x) = \|x\| \neq 0$. Противоречие.
\end{proof}

\begin{corollary}[Двойственность нормы]\label{hbcorollary4}
	Для любого $x \in E$:
	\[ \|x\| = \sup_{\|f\|_{E^*} \leqslant 1} |f(x)|. \]
\end{corollary}
\begin{proof}
	Неравенство $\leqslant$ очевидно: $|f(x)| \leqslant \|f\|\|x\| \leqslant \|x\|$.
	Равенство достигается на функционале из следствия \ref{cor:HB_norm} ($f(x)=\|x\|, \|f\|=1$).
\end{proof}\newpage
	\hypertarget{bilet19}{}\setupnav{18}{} % --- Содержимое файла materials/6_3.tex ---

\subsection{Пространства $D(\mathbb{R})$ и $D'(\mathbb{R})$. Действия над обобщёнными функциями; $\delta$-функция и $\delta$-образные последовательности. Уравнение $y' = 0$ в $D'(\mathbb{R})$.}

\begin{idea}[Смысл пространства $D'$]
	Обычные функции определяют значения в точках. Обобщенные функции определяют результат "взаимодействия" с гладким тестовым сигналом. 
	Это позволяет дифференцировать разрывные функции и описывать точечные воздействия ($\delta$-функция), которые физически существуют, но математически не являются классическими функциями.
\end{idea}

\begin{definition}[Пространство пробных функций $D$]
	Пространством \textbf{пробных (основных) функций} $D(\mathbb{R})$ называется пространство бесконечно дифференцируемых функций с компактным носителем $C_0^\infty(\mathbb{R})$.
	\begin{itemize}
		\item \textbf{Носитель} $\text{supp}(\varphi) = \overline{\{x \in \mathbb{R} \mid \varphi(x) \neq 0\}}$.
		\item Функция \textbf{финитна}, если её носитель --- компакт (отрезок).
	\end{itemize}
\end{definition}

\begin{example}[Функция-«шапочка»]
	Типичный представитель $D(\mathbb{R})$ --- функция $\omega(x, a)$:
	\[ \omega(x, a) = \begin{cases} \exp\left(-\frac{a^2}{a^2 - x^2}\right), & |x| < a \\ 0, & |x| \geqslant a \end{cases} \]
\end{example}

\begin{definition}[Сходимость в $D$]
	Последовательность $\varphi_n \to \varphi$ в $D$, если:
	\begin{enumerate}
		\item $\exists K \subset \mathbb{R}$ (компакт), такой что $\text{supp}(\varphi_n) \subseteq K$ для всех $n$;
		\item $\forall k \ge 0: \varphi_n^{(k)} \rightrightarrows \varphi^{(k)}$ (равномерная сходимость всех производных на $K$).
	\end{enumerate}
\end{definition}

\begin{proposition}[Неметризуемость]
	Топология в пространстве $D$ неметризуема.
\end{proposition}
\begin{proof}
	Предположим противное: существует метрика. Рассмотрим $\varphi_n(x) = \frac{1}{n}\omega(x, n)$. 
	Поточечно $\varphi_n \to 0$. В метрическом пространстве это порождало бы сходимость подпоследовательности к нулю. 
	Однако у любой сходящейся в $D$ последовательности должен быть \textbf{общий компактный носитель}. Но носители $[-n, n]$ расширяются на всю прямую $\implies$ противоречие.
\end{proof}

\begin{definition}[Обобщенная функция]
	\textbf{Обобщенными функциями} (распределениями) называются линейные непрерывные функционалы над пространством $D$.
	\begin{itemize}
		\item \textbf{Регулярные:} Порождаются $f(x) \in L_{1, loc}$ по правилу: $(f, \varphi) = \int_{-\infty}^{+\infty} f(x)\varphi(x) dx$.
		\item \textbf{Сингулярные:} Функционалы, не представимые в виде такого интеграла.
	\end{itemize}
\end{definition}

\begin{definition}[Дельта-функция]
	$\delta$-функция Дирака --- это сингулярный функционал: $(\delta, \varphi) = \varphi(0)$.
\end{definition}

\begin{proof}[Доказательство сингулярности $\delta$]
	Возьмем $\omega(x, a)$. $(\delta, \omega) = e^{-1}$. 
	Если $\delta$ регулярна ($f$), то по теореме Лебега при $a \to 0$ интеграл $\int f(x)\omega(x, a) dx \to 0$. 
	Но $e^{-1} \neq 0 \implies$ противоречие.
\end{proof}

\begin{proposition}[Действия над обобщенными функциями]
	Операции определяются через сопряжение:
	\begin{enumerate}
		\item \textbf{Умножение на $\psi \in C^\infty$:} $(\psi f, \varphi) = (f, \psi \varphi)$.
		\item \textbf{Дифференцирование:} $(f', \varphi) = -(f, \varphi')$. 
		\textit{Любая обобщенная функция бесконечно дифференцируема.}
		\item \textbf{Свертка:} Свойство $f * \delta = f$.
	\end{enumerate}
\end{proposition}

\begin{theorem}[Уравнение $y' = 0$]
	Решением уравнения $y' = 0$ в $D'(\mathbb{R})$ является только константа $y = C$.
\end{theorem}

\begin{proof}
	Пусть
	\[
	y' = 0 \quad \text{в } D'(\mathbb{R}).
	\]
	По определению производной распределения это означает, что для любой $\varphi \in D(\mathbb{R})$
	\[
	(y',\varphi) = -(y,\varphi') = 0.
	\]
	Следовательно,
	\[
	(y,\varphi') = 0 \qquad \forall \varphi \in D(\mathbb{R}).
	\]
	Итак, распределение $y$ обращается в нуль на всех функциях, которые являются производными пробных функций.
	
	\medskip
	\textbf{Шаг 1. Опишем множество всех производных функций из $D(\mathbb{R})$.}
	
	Обозначим
	\[
	K := \left\{\psi \in D(\mathbb{R}) \;\middle|\; \int_{-\infty}^{+\infty} \psi(x)\,dx = 0 \right\}.
	\]
	Покажем, что
	\[
	\{\varphi' \mid \varphi \in D(\mathbb{R})\} = K.
	\]
	
	\smallskip
	\textit{(i) Включение $\subseteq$.}
	Если $\psi = \varphi'$ для некоторой $\varphi \in D(\mathbb{R})$, то
	\[
	\int_{-\infty}^{+\infty} \psi(x)\,dx
	=
	\int_{-\infty}^{+\infty} \varphi'(x)\,dx
	=
	\varphi(+\infty)-\varphi(-\infty)=0,
	\]
	поскольку $\varphi$ имеет компактный носитель. Значит, $\psi \in K$.
	
	\smallskip
	\textit{(ii) Включение $\supseteq$.}
	Пусть теперь $\psi \in K$, то есть $\psi \in D(\mathbb{R})$ и
	\[
	\int_{-\infty}^{+\infty} \psi(x)\,dx = 0.
	\]
	Определим функцию
	\[
	\varphi(x) := \int_{-\infty}^{x} \psi(t)\,dt.
	\]
	Тогда $\varphi \in C^\infty(\mathbb{R})$ и
	\[
	\varphi'(x)=\psi(x).
	\]
	Осталось проверить, что $\varphi$ имеет компактный носитель.
	
	Так как $\psi$ финитна, существует отрезок $[a,b]$, вне которого $\psi=0$.
	Тогда:
	\begin{itemize}
		\item при $x<a$ имеем $\varphi(x)=0$, так как интеграл берётся по области, где $\psi=0$;
		\item при $x>b$ имеем
		\[
		\varphi(x)=\int_{-\infty}^{x}\psi(t)\,dt
		=
		\int_{-\infty}^{+\infty}\psi(t)\,dt
		=
		0.
		\]
	\end{itemize}
	Значит, $\varphi(x)=0$ вне $[a,b]$, то есть $\varphi \in D(\mathbb{R})$.
	
	Итак, действительно
	\[
	\{\varphi' \mid \varphi \in D(\mathbb{R})\}=K.
	\]
	
	\medskip
	\textbf{Шаг 2. Условие $y'=0$ означает, что $y$ зануляется на $K$.}
	
	Из предыдущего шага и равенства
	\[
	(y,\varphi')=0 \qquad \forall \varphi \in D(\mathbb{R})
	\]
	получаем:
	\[
	(y,\psi)=0 \qquad \forall \psi \in K.
	\]
	
	\medskip
	\textbf{Шаг 3. Разложим произвольную пробную функцию на сумму элемента из $K$ и фиксированной функции с интегралом $1$.}
	
	Выберем функцию $\eta \in D(\mathbb{R})$ такую, что
	\[
	\int_{-\infty}^{+\infty}\eta(x)\,dx = 1.
	\]
	(Например, можно взять подходящую нормированную ``шапочку''.)
	
	Для любой $\varphi \in D(\mathbb{R})$ положим
	\[
	\psi := \varphi - \left(\int_{-\infty}^{+\infty}\varphi(x)\,dx\right)\eta.
	\]
	Тогда
	\[
	\int_{-\infty}^{+\infty}\psi(x)\,dx
	=
	\int \varphi(x)\,dx
	-
	\left(\int \varphi(x)\,dx\right)\int \eta(x)\,dx
	=
	0.
	\]
	Следовательно, $\psi \in K$, и потому
	\[
	\varphi
	=
	\psi + \left(\int_{-\infty}^{+\infty}\varphi(x)\,dx\right)\eta.
	\]
	
	\medskip
	\textbf{Шаг 4. Вычислим действие $y$ на произвольной функции.}
	
	Так как $\psi \in K$, то $(y,\psi)=0$. Поэтому
	\[
	(y,\varphi)
	=
	(y,\psi)
	+
	\left(\int_{-\infty}^{+\infty}\varphi(x)\,dx\right)(y,\eta)
	=
	\left(\int_{-\infty}^{+\infty}\varphi(x)\,dx\right)(y,\eta).
	\]
	Обозначим
	\[
	C := (y,\eta).
	\]
	Тогда для любой $\varphi \in D(\mathbb{R})$
	\[
	(y,\varphi)
	=
	C \int_{-\infty}^{+\infty}\varphi(x)\,dx.
	\]
	
	Но правая часть есть действие регулярного распределения, порождённого постоянной функцией $C$, поскольку
	\[
	(C,\varphi)=\int_{-\infty}^{+\infty} C\,\varphi(x)\,dx
	=
	C \int_{-\infty}^{+\infty}\varphi(x)\,dx.
	\]
	Следовательно,
	\[
	y=C.
	\]
	
	Таким образом, все решения уравнения $y'=0$ в $D'(\mathbb{R})$ и только они --- постоянные распределения.
\end{proof}


\newpage % Последний
	
\end{document}